% concatenated from main.tex
\documentclass[twoside]{article}

\usepackage[accepted]{aistats2025}% If your paper is accepted, change the options for the package
% aistats2024 as follows:
%
%
% This option will print headings for the title of your paper and
% headings for the authors names, plus a copyright note at the end of
% the first column of the first page.

% If you set papersize explicitly, activate the following three lines:
%\special{papersize = 8.5in, 11in}
%\setlength{\pdfpageheight}{11in}
%\setlength{\pdfpagewidth}{8.5in}

% If you use natbib package, activate the following three lines:
%\usepackage[round]{natbib}
%\renewcommand{\bibname}{References}
%\renewcommand{\bibsection}{\subsubsection*{\bibname}}

% If you use BibTeX in apalike style, activate the following line:
%\bibliographystyle{apalike}



\usepackage{amssymb,amsfonts,amsmath,amsthm} %ams
\usepackage{enumerate,enumitem,tikz,graphicx,mathrsfs,eucal,verbatim, bbm, derivative}
\usepackage{hyperref}
\usepackage{tkz-graph, tikz-cd, svg}
\usepackage{graphicx}
\usepackage{chessfss}
 \usepackage{relsize}
\usepackage[makeroom]{cancel}
\usepackage{adjustbox}
\usepackage{diagbox}
\usepackage[thinc]{esdiff}

\usepackage{caption}
\usepackage{subcaption}


\let\svthefootnote\thefootnote
\newcommand\freefootnote[1]{%
  \let\thefootnote\relax%
  \footnotetext{\hspace{-1.5em}#1}%
  \let\thefootnote\svthefootnote%
}

 \usepackage[
    backend=biber,
    style=authoryear,
    maxcitenames=1,
    uniquelist=false
  ]{biblatex}
% \bibliographystyle{apalike}
\addbibresource{biblio.bib}



\theoremstyle{plain}% default
\newtheorem{thm}{Theorem}[section]
\newtheorem{qst}[thm]{Question}
\newtheorem{infthm}[thm]{Informal Theorem}

\newtheorem{exmp}[thm]{Example}
\newtheorem{prop}[thm]{Proposition}
\newtheorem{conj}[thm]{Conjecture}
\newtheorem{lem}[thm]{Lemma}
\newtheorem{cor}[thm]{Corollary}
\theoremstyle{definition}
\newtheorem{defn}{Definition}

\newtheorem{conjecture}{Conjecture}
\newtheorem{assm}{Assumption}
\newtheorem{prob}{Problem}
\theoremstyle{remark}
\newtheorem{rmk}{Remark}[section]


\newcommand{\samp}{\bar{x}}
\newcommand{\symm}[2]{\textnormal{Sym}_{#1}(\mathbb{R}^{#2})}
\newcommand{\symmm}[2]{\textnormal{Sym}_{#1}(#2)}
\newcommand{\edd}{\textnormal{gED}}
\newcommand{\conv}[1]{\star_{#1}}

\newcommand{\vs}[1]{\textcolor{blue}{{#1}}}
\newcommand{\gm}[1]{\textcolor{green}{{#1}}}



\begin{document}

% If your paper is accepted and the title of your paper is very long,
% the style will print as headings an error message. Use the following
% command to supply a shorter title of your paper so that it can be
% used as headings.
%
\runningtitle{Polynomial Convolutional Networks}

% If your paper is accepted and the number of authors is large, the
% style will print as headings an error message. Use the following
% command to supply a shorter version of the authors names so that
% they can be used as headings (for example, use only the surnames)
%
%\runningauthor{Surname 1, Surname 2, Surname 3, ...., Surname n}

\twocolumn[

\aistatstitle{On the Geometry and Optimization \\of Polynomial Convolutional Networks}

\aistatsauthor{Vahid Shahverdi * \And Giovanni Luca Marchetti * \And Kathl\'en Kohn *}

\aistatsaddress{KTH Royal Institute of Technology, Stockholm, Sweden} ]

\begin{abstract}
We study convolutional neural networks with monomial activation functions. Specifically, we prove that their parameterization map is regular and is an isomorphism almost everywhere, up to rescaling the filters. By leveraging on tools from algebraic geometry, we explore the geometric properties of the image in function space of this map -- typically referred to as neuromanifold. In particular, we compute the dimension and the degree of the neuromanifold, which measure the expressivity of the model, and describe its singularities. Moreover, for a generic large dataset, we derive an explicit formula that quantifies the number of critical points arising in the optimization of a regression loss.    



%We study the geometry of function spaces defined by convolutional neural networks with monomial activations. Specifically, we prove that the parametrization of such networks, once projectified, is finite birational and, in particular, it is an isomorphism almost everywhere. Consequently, we derive the dimension of the function space. Moreover, we consider the optimization of the regression loss and provide an explicit expression for the number of its critical points for a generic large dataset. Our methods are grounded in algebraic geometry and, more specifically, in the properties of Segre-Veronese varieties, to which these function spaces are closely related. 
\end{abstract}

\section{INTRODUCTION}\label{sec:intro}
Deep neural networks -- and, in general, parametric machine learning models -- define a space of functions as their parameters vary. These spaces are often referred to as \emph{neuromanifolds} \parencite{marchetti2025invitation, kohn2024geometry, calin2020neuromanifolds}. Understanding the geometry of neuromanifolds is a subtle yet fundamental challenge due to its intimate connection to the training process. Namely, neural networks learn by following a gradient flow that attracts the model to (an estimate of) the ground-truth function, which can be interpreted as minimizing a functional distance over the neuromanifold. Therefore, geometric problems over neuromanifolds -- such as nearest point problems -- are related to the learning dynamics of the corresponding models. 

\begin{figure}[t!]
    \centering
    \includegraphics[width=1.\linewidth]{Figures/segre2.png}
    \caption{Illustration of a Segre--Veronese variety parametrizing CNNs.}
    \label{fig:segre}
\end{figure}

For neural networks with polynomial activation functions, the neuromanifold is a (semi-) algebraic set, i.e., it is defined by polynomial equalities and inequalities. This enables to exploit tools from the rich field of \emph{algebraic geometry} for analyzing neuromanifolds. In particular, various algebro-geometric invariants can provide insights into fundamental machine learning aspects. \freefootnote{\quad *Equal contribution.}For example, the \emph{dimension} of the neuromanifold  plays a central role in statistical learning theory\footnote{In this context, neuromanifolds are referred to as `hypothesis spaces', and are usually considered in a combinatorial version.} since, according to the Fundamental Theorem of Learning \parencite{shalev2014understanding}, it controls the sample complexity of learnability. Moreover, the \emph{degree} is a curvature-based invariant that controls the concentration of measure around the neuromanifold and, in particular, the approximation error of the given model \parencite{basu2023hausdorff}. Dimension and degree are, in other words, basic measures of expressivity. Lastly, the \emph{Euclidean distance degree} \parencite{draisma2016euclidean} is a modern invariant counting the singularities of the distance function over the neuromanifold from an external point. In particular, it upper-bounds the number of (locally-) nearest points. Since these invariants are well-understood for a wide class of algebraic varieties, they can provide insights into the learning process of networks of various architectures. 



In this work, we consider deep \emph{Convolutional} Neural Networks (CNNs --  \cite{fukushima1979neural, lecun1995convolutional}) with monomial activation functions, and study their neuromanifolds. Our motivation is twofold. First, CNNs are popular models deployed in various signal processing domains. Historically, they have played a central role in several modern breakthroughs across deep learning \parencite{krizhevsky2012imagenet, van2016wavenet}, and have recently seen a resurgence in computer vision \parencite{liu2022convnet}. Second, the convolutional architecture is particularly suitable for algebraic geometry. Indeed, the convolution is a well-behaved algebraic operation, which translates into geometric properties for the associated neuromanifold and its parametrization. 

\subsection{Summary of Results}\label{sec:summres}
We provide theoretical results on both the geometry and the optimization of polynomial CNNs. Our arguments involve tools and ideas from algebraic geometry. In order to analyze neuromanifolds of polynomial CNNs, we study their parametrization. In this context, our main result states, informally, the following. % The main result on the geometry of neuromanifolds of polynomial CNNs establishes, informally, the following.
\begin{infthm}[Section \ref{sec:geome}]\label{thm:inf1}
Up to rescaling each filter, the parameterization of the neuromanifold of a polynomial CNN is regular and is an isomorphism almost everywhere. The remaining fibers are finite. 
\end{infthm}
Intuitively, this shows that CNNs are parametrized in an `optimal way' since, once the scaling symmetries are factored out, the parameters are in a one-to-one regular correspondence (almost everywhere) with the associated functions. This is in contrast with other neural architectures -- such as fully-connected networks -- where the parametrization exhibits large fibers and critical points, translating into redundancy in the parameter space \parencite{kileel2019expressive}. 

An immediate consequence of Informal Theorem \ref{thm:inf1} is an expression for the dimension and the degree of the neuromanifold -- see Corollary \ref{cor:finiteparam}. As mentioned in Section \ref{sec:intro}, these two invaraints play a central role in learning theory. In particular, we observe that the dimension grows linearly w.r.t. the number of layers, while the degree grows (super-) exponentially. This provides an intuitive explanation to the importance of depth in neural networks; neuromanifolds of deeper CNNs remain low-dimensional -- translating into a low sample complexity of learning -- while efficiently filling their ambient function space -- translating into expressive power.  

Informal Theorem \ref{thm:inf1} implies other geometric properties. First, it follows that the neuromanifold is closed in the Zariski topology of its ambient space (see Proposition \ref{prop:zarboundary}), and therefore in the Euclidean one as well. In terms of learning dynamics, this shows that the neuromanifold contains its asymptotic limits. Another consequence is the description of \emph{singularities} of the neuromanifold, i.e., special CNNs where the neuromanifold looks degenerate. Studying singularities of neuromanifolds and their relation to the learning dynamics is the focus of Singular Learning Theory \parencite{watanabe2009algebraic,amari2006singularities}. It follows from Informal Theorem \ref{thm:inf1} that the singular points are the simplest possible -- specifically `nodal' singularities -- which, as discussed below, translates into advantages for the optimization dynamics. Such singularities can only come from `subnetworks', meaning that they arise when (specific) weights vanish -- see Corollary \ref{cor:singpts}. The relation between singularities and subnetworks might be interpreted as a geometric explanation of the implicit bias observed in neural networks to converge towards simpler architectures during training \parencite{frankle2018lottery}.

In order to prove Informal Theorem \ref{thm:inf1}, we first remove the scaling symmetries by working in the projective space, and then to rely on tools from projective geometry. In particular, we leverage on birational geometry which, roughly speaking, concerns with properties that are preserved almost everywhere. Moreover, we factorize the projectified parametrization into a linear projection and the \emph{Segre--Veronese} embedding. The latter is a classical map in algebraic geometry whose properties are well-understood -- see Figure \ref{fig:segre} for an illustration. 

As anticipated before, understanding the geometry of the neuromanifold enables us to study the learning dynamics of polynomial CNNs. In particular,  considering the square-error regression loss, we prove the following. 
\begin{infthm}[Section \ref{sec:optim}]
For large generic datasets, the square-error loss can be rephrased as a distance minimization over the neuromanifold from an external polynomial function. There is an explicit upper bound for the number of (complex) critical points of the loss. 
\end{infthm}
The formula above is obtained via the theory of the Euclidean distance degree, which has a known expression for the Segre--Veronese variety.
% Surprisingly, for almost all datasets, the number of critical points over the neuromanifold is finite and does not depend on the dataset nor on its size (assuming it is large enough). However, the theory requires taking complex solutions into account. 
Our formula upper-bounds the number of both real and complex critical points. Since local minima are critical points, we obtain an upper bound for their number as well. Finally, we argue that  counting critical points over the neuromanifold is (almost) equivalent to counting them, up to scaling, over the parameter space, which is where optimization is performed in practice. Specifically, since the parametrization is regular (Theorem \ref{thm:inf1}), the dynamics are equivalent when considered in parameter space or over the neuromanifold. Lastly, singular points of the neuromanifold require special care; to this end, we show that they are not critical for the distance function (except for the 0-function), which is possible via our description of singularities. 

\section{RELATED WORK}
Since our work is concerned with neuromanifolds of polynomial CNNs, we review the literature around neuromanifolds and polynomial neural networks. 

\paragraph{Algebraic Geometry of Neuromanifolds.} Neuromanifolds have been analyzed via algebraic geometry in several instances. Fully-connected polynomial networks have been discussed in \parencite{kileel2019expressive, finkel2024activation}, with a focus on problems regarding dimensionality, while linear CNNs have been discussed in \parencite{kohn2022geometry, kohn2023function, shahverdi2024algebraic}, with a focus on singularities and critical points of the loss function. More recently, neuromanifolds of linear self-attention networks have been considered \parencite{henry2024geometrylightningselfattentionidentifiability}. Here, we contribute to this line of research by discussing the geometry of neuromanifolds defined by polynomial CNNs. While the previous literature focuses mainly on the linear case -- with only partial results for the general polynomial one -- our work is the first to provide a comprehensive description of both the parametrization and the neuromanifold of a polynomial model. 
% \paragraph{Dynamics of Neural Gradient Descent.} Several works have analysed the gradient flow of the loss function of a neural network, which lies at the heart of the training process of contemporary machine learning models. ... \emph{hidden symmetries} ...
\paragraph{Polynomial Activation Functions.} Standard activation functions for neural networks -- such as the Rectified Linear Unit (ReLU) -- are not polynomial. Nonetheless, networks with polynomial activations have been considered in the literature. Such networks are universal interpolators and approximators \parencite{constantinescu2023interpolation}. Various flavors of polynomial activations have been discussed, ranging from monomial versions of ReLU \parencite{berradi2018symmetric, li2019powernet}, to rational functions \parencite{boulle2020rational, telgarsky2017neural}, to piece-wise polynomial functions \parencite{lopez2019piecewise, hou2017convnets}. Finally, quadratic activations have appeared in neuroscience for modelling biological neural networks \parencite{adelson1985spatiotemporal}, and have sometimes been considered in theoretical studies of deep learning phenomena \parencite{gromov2023grokking, marchetti2024harmonics}. 

\section{BACKGROUND}\label{sec:back}
In this section, we introduce the necessary background around convolutional neural networks, as well as the relevant tools from algebraic geometry. 

\subsection{Polynomial Convolutional Networks}\label{sec:polconvs}
We start by introducing convolutional neural networks. For simplicity of notation, we focus on one-dimensional convolutions -- all theory and results can be extended verbatim to the higher-dimensional case. Given integers $k, s, d' \in \mathbb{N}$ representing filter size, stride, and output dimension respectively, the convolution between a filter $w \in \mathbb{R}^k$ and an input vector $x \in \mathbb{R}^d$, with $d = s(d' -  1) + k$, is the vector $ w \conv{s} x \in \mathbb{R}^{d'}$ defined for $ 0 \leq i < d'$ as: 
\begin{equation}
\label{eq:conv_star}
 (w \conv{s} x)[i] =   \sum_{0 \leq j < k} w[j] \  x[si + j].
\end{equation}
The convolution in Equation \eqref{eq:conv_star} is linear in $x$ and is represented by a $d' \times d$  \emph{Toeplitz matrix}. For example,  given $k=3$, $s=2$, and $d'=3$, the corresponding matrix is 
\[\begin{pmatrix}
    w[0]&w[1]&w[2]&0&0&0&0\\
    0&0&w[0]&w[1]&w[2]&0&0\\
    0&0&0&0&w[0]&w[1]&w[2]
\end{pmatrix}.\]

Moreover, the composition of convolutions can be rephrased as polynomial multiplications. We associate to the filter the following homogeneous bivariate polynomial of degree $k-1$ in $(a^s,b^s)$: 
\begin{equation}
    \label{eq:poly_correspondence}
    \pi_s(w) = \sum_{0 \leq i < k}w[i] \ a^{s(k-i-1)}b^{si} .
\end{equation}
Then the following property holds: an iterated convolution $v \conv{s} (w \conv{t} x)$ with respective strides $s, t$ coincides with a convolution $q \conv{st} x$ whose associated polynomial satisfies $\pi_{1}(q) = \pi_{t}(v) \ \pi_{1}(w)$.  

Convolutions can be composed in order to construct deep convolutional networks. To this end, fix an integer $L \in \mathbb{N}$ and sequences $\mathbf{k}, \mathbf{s} \in \mathbb{N}^{L}$, $\mathbf{d} \in \mathbb{N}^{L+1}$ such that $d_i = s_i(d_{i+1} + k_i)$ for all $i$. Moreover, consider a map $\sigma \colon \mathbb{R} \rightarrow \mathbb{R}$ playing the role of the activation function. 
\begin{defn}
A \emph{Convolutional Neural Network} (CNN) with weights $\mathbf{w} = (w_{0}, \ldots, w_{L-1}) \in \bigoplus_{0 \leq i < L}\mathbb{R}^{k_i}$ is the map $\varphi_\mathbf{w} \colon \mathbb{R}^{d_0} \rightarrow \mathbb{R}^{d_L}$ given by:
\begin{equation}
% \begin{array}{rccl}
% \varphi_\mathbf{w} \colon & \mathbb{R}^{d_0}  & \rightarrow & \mathbb{R}^{d_L} \\
%  & x & \rightarrow & w^{L-1} \conv{s_{L-1}} \sigma \left( w^{L-2} \conv{s_{L-2}} \cdots \star  \sigma ( w^0 \conv{s_{0}} x   ) \right),
% \end{array}
\varphi_\mathbf{w}(x) = w_{L-1} \conv{s_{L-1}} \sigma \left(\cdots \conv{s_1}  \sigma ( w_0 \conv{s_{0}} x   ) \right),
\end{equation}
where $\sigma$ is applied coordinate-wise.
\end{defn}
Given $r \in \mathbb{N}$, denote by $\sigma_r(x) = x^r$ the power function. CNNs with activation function $\sigma = \sigma_r$ are called \emph{polynomial}, and they are called \emph{linear} if $r = 1$ i.e., if $\sigma$ is the identity function. Polynomial CNNs define homogeneous polynomial functions of degree $r^{L-1}$, meaning that $\varphi_\mathbf{w} \in \symm{r^{L-1}}{d_0}^{d_L}$. Here, $\symm{\alpha}{\beta}$ denotes the space of symmetric tensors of degree $\alpha$ in $\beta$ variables, i.e., the vector space of dimension $\binom{\beta + \alpha - 1}{\alpha}$ whose canonical basis is given by monomials of degree $\alpha$ in $x[0], \ldots, x[\beta -1]$. 
\begin{defn}
The \emph{neuromanifold} of a polynomial CNN is the image in $\symm{r^{L-1}}{d_0}^{d_L}$ of the parametrization, i.e.,
\begin{equation}
\mathcal{M}_{\mathbf{d}, \mathbf{k}, \mathbf{s}, r} = \left\{\varphi_\mathbf{w} \ | \  \mathbf{w} \in \mathbb{R}^{|\mathbf{k}|} \right\}, 
\end{equation}
\end{defn}
where $|\mathbf{k}| = k_0 + \cdots + k_{L-1}$. By the Tarski-Seidenberg theorem, the neuromanifold is a semi-algebraic set, i.e., it can be defined by polynomial equalities and inequalities. However, we will show that it is actually an algebraic \emph{variety}, meaning that it is closed in Zariski topology of $\symm{r^{L-1}}{d_0}^{d_L}$, or, equivalently, that it can be defined by polynomial equalities alone -- see Proposition \ref{prop:zarboundary}. Lastly, throughout this work we will often consider complex coefficients. To this end, we denote by $\mathcal{M}_{\mathbf{d}, \mathbf{k}, \mathbf{s}, r}^\mathbb{C}$ the complexification of the neuromanifold. The latter can be defined by replacing real numbers $\mathbb{R}$ with complex ones $\mathbb{C}$ in all  definitions of this section. 

\subsection{Segre--Veronese Varieties}
Here, we recall the notions of projective spaces and Segre--Veronese varieties, together with their basic properties. As we will show, these are closely related to the neuromanifolds of CNNs. To this end, fix a field $\mathbb{K}$ (e.g., $\mathbb{R}$ or $\mathbb{C}$). For a vector space $V$ over $\mathbb{K}$, its \emph{projectification} $\mathbb{P} V$ consists of all the non-vanishing vectors up to rescaling or, equivalently, of all the lines
in $V$. Formally, it can be defined as the quotient $\mathbb{P} V =(V \setminus \{ 0\}) / (\mathbb{K} \setminus \{ 0\})$, where $\mathbb{K}$ acts on $V$ via scalar multiplication. Projective spaces will be relevant for us since, as we shall see, it is convenient to consider the parameters of polynomial CNNs -- i.e., the filters -- up to rescaling. 

Now, fix $k \in \mathbb{N}$, let $\mathbf{m}, \mathbf{p} \in \mathbb{N}^k$, and consider $k$ vector spaces $V_1, \ldots, V_k$ over $\mathbb{K}$, with $\dim(V_i) = p_i + 1$. 

\begin{defn}\label{def:segveremb}
    The \emph{Segre--Veronese embedding} is the map
    \begin{equation}
        \nu_{\mathbf{m}, \mathbf{p}} \colon \prod_{1 \leq i \leq k}\mathbb{P}V_i \rightarrow \mathbb{P} \bigotimes_{1 \leq i \leq k} 
      \symmm{m_i}{V_i}  
    \end{equation}
defined by taking tensor products of symmetric powers of vectors in the corresponding spaces. The \emph{Segre--Veronese variety} $\mathcal{V}_{\mathbf{m}, \mathbf{p}}$ is the image of $ \nu_{\mathbf{m}, \mathbf{p}}$.
\end{defn}
Explicitly, if $V_i = \mathbb{R}^{p_i + 1}$, then the Segre--Veronese variety consists of (tensor) products of $k$ monomials of corresponding degree $m_i$. In particular, $ \nu_{\mathbf{m}, \mathbf{p}} $ is an embedding, as suggested by the nomenclature, and $\mathcal{V}_{\mathbf{m}, \mathbf{p}}$ is a smooth projective variety of dimension $| \mathbf{p}| = p_1 + \cdots +  p_{k}$. When $k=1$, we refer to $\nu_{m,p}$ simply as Veronese embedding. For example, for $k=2$ and $p_1 = p_2 = m_1 = m_2 = 1$, $\mathcal{V}_{\mathbf{m}, \mathbf{p}}$ coincides with a smooth quadric in the three-dimensional projective space $\mathbb{P}(V_1 \otimes V_2)$, which can be represented as a hyperboloid -- see Figure \ref{fig:segre}. As we shall see, since $\varphi_\mathbf{w}(x)$ depends polynomially on $\mathbf{w}$, the prametrization and the neuromanifold of polynomial CNNs are closely related to the Segre-Veronese embedding and variety, respectively. 

\subsection{Euclidean Distance Degree}\label{sec:eucdd}
We now recall the \emph{Euclidean distance degree} \parencite{draisma2016euclidean} -- an invariant in algebraic geometry that will be central in our work. Intuitively, this invariant counts the number of critical points of the Euclidean distance function from (the smooth locus of) an algebraic variety to a fixed external point. Formally, let $X \subset \mathbb{R}^n$ be an algebraic variety, whose smooth locus is denoted by $X_{\textnormal{reg}}$. Let $A$ be a $n \times n$ symmetric positive-definite matrix inducing a distance $\textnormal{d}_A(x, y)^2 = (x-y)^\top A \ (x-y) $ for $x, y \in \mathbb{R}^n$. We consider the (squared) distance function from an anchor $u \in \mathbb{R}^n \setminus X$ to the smooth locus of $X$: 
\begin{equation}\label{eq:distfunc}
\begin{array}{ccc}
  X_{\textnormal{reg}}  & \rightarrow & \mathbb{R}_{\geq 0} \\
  x & \mapsto & \textnormal{d}_A(x, u)^2.
\end{array}
\end{equation}
For what follows, it is necessary to consider complex coefficients. To this end, we complexify the above map, extending it as $X^\mathbb{C}_\textnormal{reg} \rightarrow \mathbb{C}$, where $X^\mathbb{C} \subset \mathbb{C}^n$ denotes the complexification of $X$. 
\begin{defn}\label{def:edd}
The \emph{Euclidean distance degree} of $X$ with respect to $A$ is the  number of complex critical points of the complexified map $X^\mathbb{C}_\textnormal{reg} \rightarrow \mathbb{C}$ for generic~$u$. 
\end{defn}
Geometrically, a point $x \in X_\textnormal{reg}$ is critical for the distance map if, and only if, $x - u$ is perpendicular to the tangent space of $X$ at $x$ according to the scalar product induced by $A$. Moreover, it is possible to extend the notion of Euclidean distance degree to projective varieties $X \subset \mathbb{P}\mathbb{R}^{n}$ by considering their de-projectification, i.e., the corresponding affine cone in $\mathbb{R}^n$. 


\begin{figure}[!ht]
    \centering
    \includegraphics[width=0.43\linewidth]{Figures/eddbezier2.png}
    \caption{Distance function from an anchor to a curve, visualized as a color gradient. The critical values are denoted by dotted lines.}
    \label{fig:eddbezier}
\end{figure}

For generic $A$, the Euclidean distance degree is the same, achieving its maximal value \parencite[Theorem 6.11]{draisma2016euclidean}. This number is referred to as \emph{generic} Euclidean distance degree of $X$, denoted by $\edd(X) \in \mathbb{N}$. The following result -- which will be necessary in our forthcoming discussion -- provides a closed formula for the generic Euclidean distance degree of Segre--Veronese varieties. 

\begin{thm}[\cite{kozhasov2023minimal}]\label{th:segverdeg}
The generic Euclidean Distance degree of the Segre--Veronese variety is:
\begin{align*}
    \edd(\mathcal{V}_{\mathbf{m}, \mathbf{p}}) = &\sum_{0 \leq i \leq |\mathbf{p}|} (-1)^i \left( 2^{|\mathbf{p}| + 1 - i} - 1 \right)(|\mathbf{p}| - i)! \\ 
    &\sum_{\substack{|\mathbf{\alpha}|=i \\ \forall j \ \alpha_j \leq p_j}} \prod_{1\leq j \leq k} \frac{\binom{p_j + 1}{\alpha_j}}{(p_j - \alpha_j)!} \ m_j^{p_j - \alpha_j},
\end{align*}
where $|\mathbf{p}| = p_1 + \cdots + p_{k}$. 
\end{thm}
As anticipated in Section \ref{sec:intro}, machine learning models are trained to minimize, ideally, the distance to the ground-truth function. Therefore, the distance function over neuromanifolds and its critical points are closely related to the learning process. This motivates the study of the Euclidean distance degree of neuromanifolds, which we will explore for polynomial CNNs in Section \ref{sec:optim}.  

\section{CONVOLUTIONAL NEUROMANIFOLDS}\label{sec:convneu}
In this section, we study the neuromanifolds of polynomial CNNs, and present the main results of this work. We first discuss geometric properties -- such as dimension, projectification, and relation to Segre--Veronese varieties -- and then proceed to compute the generic Euclidean distance degree of the neuromanifold. We provide most of the proofs in the appendix -- see Section \ref{sec:remproof}. Moreover, in Section \ref{sec:toyexmp}, we provide a toy example where all the quantities discussed here can be computed explicitly.
\begin{figure*}[th!]
\captionsetup[subfigure]{justification=centering}
    \centering
    \hspace{1em}
    \begin{subfigure}[b]{.27\linewidth}
        \centering
        \includegraphics[width=\linewidth]{Figures/r_2_resized.png}
        \subcaption*{$r = 2$}
    \end{subfigure}
    \hfill
    \begin{subfigure}[b]{.27\linewidth}
        \centering
        \includegraphics[width=\linewidth]{Figures/r_5_resized.png}
        \subcaption*{$r = 5$}
    \end{subfigure}
    \hfill
    \begin{subfigure}[b]{.27\linewidth}
        \centering
        \includegraphics[width=\linewidth]{Figures/r_10_resized.png}
        \subcaption*{$r = 10$}
    \end{subfigure}
    \caption{Visualization of two-dimensional charts of neuromanifolds (over $\mathbb{R}$) corresponding to $(k_0,k_1)=(2,2)$ projected orthogonally to $\mathbb{R}^3$, with varying activation degree. }
    \label{fig:neur}
\end{figure*}
\subsection{Geometry}\label{sec:geome}
Throughout this section, we will leverage on principles and tools from algebraic geometry. Therefore, it will be convenient to work with complex coefficients. For simplicity, we will abuse the notation and use the symbols from Section \ref{sec:polconvs} to denote the analogous complex object, obtained by replacing $\mathbb{R}$ with $\mathbb{C}$ in the corresponding definition. While all the results of this section are stated and proved over complex numbers, almost all of them extend a posteriori to the real case -- see Remark \ref{rem:realcase}.

We start with a simple but powerful property of convolutions.
\begin{lem}\label{lem:fullrk}
    For $w \in \mathbb{C}^k \setminus \{ 0 \}$ and $s \in \mathbb{N}$, the convolution map $x \mapsto w \conv{s} x$ has full rank. 
\end{lem}
\begin{proof}
See Section \ref{app:fullrk} in the appendix.  
\end{proof}
Consider now a polynomial CNN  $\varphi_\mathbf{w}$ parametrized by weights $\mathbf{w} \in \mathbb{C}^{|\mathbf{k}|}$. 

% \begin{exmp}
%     Consider a one-dimensional PCNN with filter sizes $\mathbf{k}=(3,2)$, strides $\mathbf{s}=(1,1)$ and activation function $\rho_2(x)$. If the input data has the symbolic form $\mathbf{x} = \begin{bmatrix}
%             x_1&x_2&x_3&x_4
%         \end{bmatrix}^\top$ and let $\tilde{\mathbf{x}}=\begin{bmatrix}
%             x_1^2&x_1x_2& x_1x_3 & x_2^2& x_2x_3& x_2x_4& x_3^2& x_3x_4&0& x_4^2
%         \end{bmatrix}^\top$, then the output of this network can be presented by 
%     \begin{align*}
%         &\begin{bmatrix}
%             w_{21}&w_{22}
%         \end{bmatrix}\rho_2 \begin{bmatrix}
%             w_{11}&w_{12}&w_{13}&0\\
%             0&w_{11}&w_{12}&w_{13}
%         \end{bmatrix}\cdot {\mathbf{x}}\\
%         =& \begin{bmatrix}
%             w_{21}&w_{22}
%         \end{bmatrix} \begin{bmatrix}
%             w_{11}^2 & 2w_{11}w_{12}& 2w_{11}w_{13} & w_{12}^2 & 2w_{12}w_{13} & 0 & w_{13}^2 &0&0&0\\
%             0&0&0& w_{11}^2 & 2w_{11}w_{12}& 2w_{11}w_{13} & w_{12}^2 & 2w_{12}w_{13} & 0 & w_{13}^2 
%         \end{bmatrix}\cdot \tilde{\mathbf{x}}.
%     \end{align*}
%     The neuromanifold of this network, denoted by $\mathcal{M}_{\mathbf{k},\mathbf{s},\rho_2}$, is a subset of the neuromanifold of an LCNN with filter sizes $\mathbf{k'} = (7, 2)$ and strides $\mathbf{s'} = (3, 1)$.
% \end{exmp}




\begin{cor}\label{cor:basepts}
    If $\varphi_\mathbf{w}(x) = 0$ for all $x \in \mathbb{C}^{d_0}$, then $w_i = 0$ for some $i$. 
\end{cor}
\begin{proof}
Suppose that $w_i \not = 0 $ for all $i$. By Lemma \ref{lem:fullrk}, for all $i$ the map $x \mapsto \sigma_r( w_i \conv{s_i} x)$ is open, i.e., it sends open sets to open ones. Openness is preserved by composition and, in particular,
$\varphi_\mathbf{w} \not = 0$, as desired. 
\end{proof}

\begin{rmk}\label{rmk:basepts}
The property established by Corollary \ref{cor:basepts} is characteristic of (polynomial) CNNs, and does not hold for other network architectures. For example, a fully-connected network with arbitrary activations and at least two layers will not satisfy it. To illustrate this, consider a $2$-layer fully-connected network $W_1 \cdot \sigma(W_0 \cdot x)$, where $x \in \mathbb{R}^{d_0}$ is the input, $W_0 \in \mathbb{R}^{d_1 \times d_0}$ and $W_1 \in \mathbb{R}^{d_2 \times d_1}$ are the weight matrices, and $\sigma \colon \mathbb{R} \rightarrow \mathbb{R}$ is the activation function. Assuming $d_1 > 1$, if the first two rows of $W_0$ coincide and the other ones vanish, while the first two columns of $W_1$ are opposite, then the network vanishes on all $x \in \mathbb{R}^{d_0}$. 
\end{rmk}



As a consequence of Corollary \ref{cor:basepts}, it is possible to projectify the neuromanifold and its parametrization. Indeed, the result implies that the map $\mathbf{w} \mapsto \varphi_\mathbf{w}$ parametrizing the neuromanifold induces an algebraic morphism between the corresponding projective spaces: 
\begin{equation}\label{eq:projsegver}
    \overline{\varphi}:\prod_{0 \leq i < L} \mathbb{P} \mathbb{C}^{k_i} \rightarrow  
    \mathbb{P} \symmm{r^{L-1}}{\mathbb{C}^{d_0}}^{d_L}.  
\end{equation} 
We denote by $\mathbb{P}\mathcal{M}_{\mathbf{d}, \mathbf{k}, \mathbf{s}, r}^\mathbb{C}$ the image of $\overline{\varphi}$ and deem it \emph{projective neuromanifold}. We deduce the following relevant property of the (projective) neuromanifold. 
\begin{prop}\label{prop:zarclosed}
Both $\mathbb{P}\mathcal{M}_{\mathbf{d}, \mathbf{k}, \mathbf{s}, r}^\mathbb{C}$ and $\mathcal{M}_{\mathbf{d}, \mathbf{k}, \mathbf{s}, r}^\mathbb{C}$ are closed in the Zariski topology of their respective ambient spaces. 
\end{prop}
\begin{proof}
See Section \ref{app:zarclosed} in the appendix. 
\end{proof}
In order to analyse $\overline{\varphi}$, we provide another technical result that enables to rephrase polynomial CNN layers in terms of convolutions. To this end, fix a filter $w \in \mathbb{C}^k$, a stride $s$, and an exponent $r$. 
% \begin{prop}\label{prop:converonese}
%  For every $x \in \mathbb{C}^n$, $\sigma_r(w \conv{s} x)$ can be written as a convolution $\tilde{w} \conv{\tilde{s}} \tilde{x}$ with stride $\tilde{s} = s \binom{r+k-2}{k-1}$ and
% filter size  $\tilde{k}=(k-1)\binom{r+k-2}{k-1}+1$. Moreover, $\tilde{w}$ and consist of monomials in $w$, while $\tilde{x}$ consists of monomials in $x$ and arbitrarily-chosen scalars. 
% \end{prop}
\begin{prop}\label{prop:converonese}
 For every $x \in \mathbb{C}^n$, $\sigma_r(w \conv{s} x)$ can be written as a convolution $\tilde{w} \conv{\tilde{s}} \tilde{x}$ with stride $\tilde{s} = s \binom{r+k-1}{r}$ and
filter size  $\tilde{k}=\binom{r+k-1}{r}$. Moreover, $\tilde{x}$ and $\tilde{w}$ consist of monomials in $x$ and $w$, respectively.  
\end{prop}
\begin{proof}   
See Section \ref{app:converonese} in the appendix. 
%Consequently, we define $\tilde{x}$ as follows. Let $A_1$ be a vector consisting of all monomials $x[0]^{a_0} \cdots x[k-1]^{a_{k-1}}$ of degree $r$ in lexicographic order such that $a_0 \geq 1$. There are exactly $\binom{r+k-2}{k-1}$ such monomials. Next, we construct $A_j$ for $1 < j < n$ by shifting all the indices in $A_{j-1}$ by one unit. If any term in a monomial has an index $\geq n$, we insert an arbitrary scalar. Finally, we set $A_n$ to a single-element vector $(x[n]^r)$. We define $\tilde{x}$ as the vector of size $(n-1) \binom{r+k-2}{k-1} + 1$ given by the concatenation of all $A_i$ in order. Moreover, we construct $\tilde{w}$ by writing the corresponding coefficients in Equation \ref{eq:newtexp} of the monomials in $\tilde{x}$. Since
%$\tilde{s}+1$ is the position of the monomial $x[s+1]^r$ in $\tilde{x}$, we have $\tilde{s} = s  \binom{r+k-2}{k-1}$. Similarly, the filter size $\tilde{k}$ is the position of $x[k]^r$ in $\tilde{x}$, which is precisely $(k-1) \binom{r+k-2}{k-1} + 1$.
\end{proof}
% \begin{rmk}
%     \label{rmk:L-layer_conv}
%     Proposition \ref{prop:converonese} implies, by induction, that for an arbitrary number of layers $L$, $\varphi_{\mathbf{w}}(x)$ can be expressed as a convolution $\tilde{w} \conv{\tilde{s}} \tilde{x}$ for some $\tilde{w}$, $\tilde{s}$, and $\tilde{x}$. Even further, by construction, it is possible to `shrink' $\tilde{x}$ by removing the repeated monomials, and adjust the filter accordingly. 
% \end{rmk}
Proposition \ref{prop:converonese} implies, inductively, that $\overline{\varphi}$ factors via the Segre--Veronese embedding (Definition \ref{def:segveremb}). Specifically, let $\mathbf{m} = (r^{L-1}, \ldots, r, 1)$ and $\mathbf{p} = (k_0 - 1, \ldots, k_{L-1} - 1)$. Then the following diagram commutes: 
\begin{equation*}
% \begin{adjustbox}{width=.48\textwidth,center}
\begin{tikzcd}
	{\displaystyle \prod_{0 \leq i < L} \mathbb{P}\mathbb{C}^{k_i} } && \mathbb{P}\mathcal{M}_{\mathbf{d}, \mathbf{k}, \mathbf{s}, r}^\mathbb{C} \\
	& \mathcal{V}_{\mathbf{m}, \mathbf{p}}
	\arrow["\overline{\varphi}",from=1-1, to=1-3]
	\arrow["{ \nu_{\mathbf{m}, \mathbf{p}}}"', from=1-1, to=2-2]
	\arrow["\Lambda"', from=2-2, to=1-3]
\end{tikzcd}
% \end{adjustbox}
\end{equation*}
Here, $\Lambda$ is (the projectification of) a linear map restricted to the Segre--Veronese variety. In the next section, this factorization  enables us to compute the neuromanifold's generic Euclidean distance degree. 


% Up to change of projective coordinates, we can assume that $\Lambda$ is a projection. Here, given vector spaces $V, V'$, a projection is the algebraic morphism $\mathbb{P}(V \oplus V') \setminus \mathbb{P}V' \rightarrow \mathbb{P}V$ induced by the map onto the direct summand $V \oplus V' \rightarrow V$. Based on this, we aim to exploit properties of projections to study $\overline{\varphi}$. To this end, 
% \begin{lem}
% \label{lem:projection_away}
%  Let $V, V'$ be finite-dimensional vector spaces and $X \subset \mathbb{P} (V \oplus V')$  be a projective variety that does not intersect $\mathbb{P}V'$. Then the restricted projection $X \rightarrow \mathbb{P}V$ is finite.
% \end{lem}
% \begin{proof}
%  By considering a flag of $V'$ -- i.e., a filtration $V_0 = 0 \subset V_1 \subset \cdots \subset V_k = V'$ with $\textnormal{dim}(V_{i+1}) = \textnormal{dim}(V_i) + 1$ -- it is possible to factor the projection into projections onto projective hyperplanes. Therefore, without loss of generality, we can assume that $\textnormal{dim}(V') = 1$, i.e., that $\mathbb{P}V' = \{ p\}$ is a point. But then the fibers of the projection are projective lines passing through $p$. By hypothesis, such lines are not contained in $X$ and, therefore, they intersect $X$ in a finite set.
%  \vs{why the projection isn't a rational map and it is a morphism? should we state it?}
%  %the fiber $\Lambda_p^{-1}(v)$ is an affine line, which is proper as well. Thus, $\#\Lambda_p^{-1}(v) < \infty$. By \parencite[III,§11, Exercise 11.2]{hartshorne2013algebraic}, a proper morphism with finite fibers is finite.
% \end{proof}


% In the linear case, i.e., when \( r = 1 \), and for strides \( s_i > 1 \) for all \( i \in \{1, \ldots, L-2\} \), the parameterization map \(\mu\) is birational \parencite[Corollary 5.4]{kohn2023function}. In the next theorem, we show that when \( r > 1 \), almost every function \(\varphi_w\) in the projectivization of the neuromanifold can be represented by exactly one tuple of filters \( w \), thus making \(\mu\) birational. This implies that the varieties \(\prod_{0 \leq i < L} \mathbb{P}^{k_i - 1}\) and \(\mathbb{P}(\mathcal{M}_{\mathbf{d}, \mathbf{k}, \mathbf{s}})\) are essentially the same in terms of their most fundamental geometric and algebraic properties.
We now provide the main results of this section. Recall that the differential of a differentiable map is defined as its local linearization, and is explicitly computable in given coordinates as the Jacobian. A differentiable map is \emph{regular} if its differential has maximal  rank at every point of the domain. 
\begin{thm}\label{thm:reg}
Assume that $r>1$. Then the projectified parametrization $\overline{\varphi}$ is regular. 
\end{thm}
\begin{proof}
The argument requires several technical computations around the differential of $\varphi$, whose details we provide in the appendix, Section \ref{sec:proofdet}. By Euler's theorem on homogeneous functions, showing that the differential of $\overline{\varphi}$ is injective at some $\mathbf{w}$ is equivalent to showing that the kernel of the differential of $\varphi$ at $\mathbf{w}$ has dimension $L-1$. The latter follows immediately from Proposition \ref{prop:kerdiff}, which provides $L-1$ independent generators for such kernel.     
\end{proof}
Next, we show that the map $\overline{\varphi}$ is \emph{birational}, meaning that it is generically one-to-one. In other words, `almost all' the fibers of $\overline{\varphi}$ are singletons. We also show that the remaining non-singleton fibers are finite.%Below, we show that the parametrization of a polynomial CNN is birational, with the remaining non-singleton fibers being finite.
\begin{thm}\label{thm:mainbir} 
Assume that $r > 1$. If $\overline{\varphi}_{\mathbf{w}} = \overline{\varphi}_{\mathbf{v}}$ for some $\mathbf{w}, \mathbf{v} \in  \prod_{i} \mathbb{P}\mathbb{C}^{k_i}  $, then $\mathbf{w}$ and $\mathbf{v}$ are related by a shift of their indices. More precisely, there exist $t_0, \ldots, t_{L-1} $ such that for all $i$:
\begin{align}
w_i &= (0, \ldots, 0, v_i[0], \ldots, v_i[k_i - t_i - 1]), \\ 
v_i &= (v_i[0], \ldots, v_i[k_i - t_i - 1], 0, \ldots, 0),
\end{align}
or vice versa. In particular,  $\overline{\varphi}$ is birational on its image, and all the fibers are finite.   
\end{thm}
\begin{proof}
% From Proposition \ref{cor:finiteparam}, we know that $\overline{\varphi}$ is finite. By \cite[Proposition 4.6.7]{grothendieck1961elements}, if a finite morphism is a closed immersion at a point, then it is a closed immersion in a (Zariski) neighborhood of that point. Intuitively, this result is an algebro-geometric analogue of the inverse function theorem in differential geometry. Therefore, in order to show that $\overline{\varphi}$ is generically one-to-one on its image, it suffices to exhibit a point at which it is a closed immersion. By \cite[Theorem 19.1.1]{Vakil2010MATH2F}, it is enough to find a tuple of filters (up to multiplicative scalar) $\mathbf{w}$ such that the fiber of $\overline{\varphi}_\mathbf{w}$ is the singleton $\{\mathbf{w}\}$ and such that $\overline{\varphi}$ is regular at at $\mathbf{w}$. But since by Theorem \ref{thm:reg} $\overline{\varphi}$ is regular at every point, it suffices to find $\mathbf{w}$ whose fiber is a singleton.  

% We claim that the tuple of filters $ \mathbf{w} = (w_0, \ldots, w_{L-1} )$ with $w_i = (1, 0, \ldots, 0)$ is as desired. The corresponding CNN is $\varphi_\mathbf{w} (x)$, which is a vector with entries of the form  $x[i]^{r^{L-1}}$. We prove that there is no other filter tuple $\mathbf{w'}$ (up to scaling) such that $\overline{\varphi}_\mathbf{w} = \overline{\varphi}_\mathbf{w'}$. By contradiction if such a $\mathbf{w'}$ exists, then by Remark \ref{rmk:identity_input}, we get
% \begin{equation*}
% \begin{aligned}
%     & W_{L-1} \left( W_{L-2} \left( \cdots \left( W_2 (W_1)^{\odot r} \right)^{\odot r} \cdots \right) \right)^{\odot r} \\
%       = \, & W'_{L-1} \left( W'_{L-2} \left( \cdots \left( W'_2 (W'_1)^{\odot r} \right)^{\odot r} \cdots \right) \right)^{\odot r}.
% \end{aligned}
% \end{equation*}
% The left-hand side of this identity is a Toeplitz matrix with filter $w$ whose corresponding polynomial, in the notation from Section \ref{sec:polconvs}, is
% \[
% \pi_1(w) = a^{(k_0 - 1) + (k_1 - 1)s_0 + \cdots + (k_{L-1} - 1)s_0 \cdots s_{L-2}}.
% \]
% Again, by using Equation \ref{eq:poly_correspondence}, we can deduce that $w'_{L-1}$ is equal to $w_{L-1}$ up to rescaling. This holds for other filters $w'_i$ as well, since $\bullet^{\odot r}$ is a point-wise power, implying that $\pi_1(w_i')$ is equal to $a^{k_i - 1}$ up to rescaling. 
The argument involves a delicate induction over the number of layers -- see Section \ref{app:mainbir} in the appendix. 
\end{proof}
Theorem \ref{thm:mainbir} can be intuitively interpreted as stating that the projective neuromanifold is isomorphic to a Segre--Veronese variety `almost everywhere'. For $r = 1$ -- i.e., for linear CNNs -- the same result has been shown by \cite{kohn2023function} with the additional assumption on strides $s_i > 1$ for all $0 \leq i \leq L - 2$. Moreover, in the same work the authors show that the parametrization of linear CNNs has critical points, in contrast to the polynomial case. The argument is different from ours and is based on the interpretation of convolutions as polynomial multiplication. 

\begin{rmk}\label{rmk:singsmall}
From the proof of Theorem \ref{thm:mainbir}, it is evident that the shifts must satisfy arithmetic constraints involving the strides. Specifically, consider the shifts $t_i$ together with their sign, which is set positive if $w_i$ has zeros on the left, and negative otherwise. Define recursively $\tilde{t}_{-1} = 0$, $\tilde{t}_i = t_i + \tilde{t}_{i-1} / s_{i-1}$ for $i \geq 0$. Then $\tilde{t}_i$ must be an integer, and moreover $\tilde{t}_{L-1}$ must vanish. By the definition of (iterated) convolutions, this implies that if we truncate each filter by removing its zeros on the right or on the left, the resulting CNN coincides with $\varphi_{\mathbf{w}}$, albeit on a restricted domain with less input variables. Put simply, the CNNs whose (projectified) fibers are not singletons are the ones that can be defined by smaller architectures, i.e., special
fibers arise from `subnetworks'.    
\end{rmk}

An immediate consequence of Theorem \ref{thm:mainbir} is an expression for the two fundamental invariants of the neuromanifold: the dimension and the degree. Recall that for a given projective/affine variety, the latter counts, roughly speaking, the number of intersections with a generic linear subspace of dimension equal to the co-dimension of the variety. 
\begin{cor}\label{cor:finiteparam}
Assume that $r>1$. The dimension and degree of the neuromanifold are: 
\begin{align}
    \textnormal{dim}(\mathcal{M}_{\mathbf{d}, \mathbf{k}, \mathbf{s}, r}^\mathbb{C}) &= |\mathbf{k}| - L + 1, \\
    \deg(\mathcal{M}_{\mathbf{d}, \mathbf{k}, \mathbf{s}, r}^\mathbb{C}) &= (|\mathbf{k}|-L)!\prod_{0 \leq j < L} \frac{r^{(L - j - 1)(k_j - 1)}}{(k_j - 1)!}.
\end{align}
\end{cor}
\begin{proof}
See Section \ref{app:finiteparam} in the appendix. 
\end{proof}
In particular, note that if all the filter sizes $k_i := k$ are equal, then the degree is $(L(k-1))! \ r^{L(L-1)(k-1) / 2 }/ (k-1)!^L$, which grows faster than $r^{L^2/2}$ w.r.t. $L$. Therefore, as anticipated in Section \ref{sec:summres}, the dimension grows linearly w.r.t. the depth of the network, while the degree grows (super-) exponentially.  

Another consequence of the main results in this section is a description of the singular points of the neuromanifold. Since $\overline{\varphi}$ is regular and finite, the singularities of $\mathbb{P}\mathcal{M}_{\mathbf{d}, \mathbf{k}, \mathbf{s}, r}^\mathbb{C}$ coincide with CNNs whose (projectified) fiber is not a singleton. Such singularities are of \emph{nodal} type, i.e., points where the neuromanifold self-intersects, creating a finite number of (potentially overlapping) tangent spaces -- see Figure \ref{fig:neur} (center and right) for an illustration. By Remark \ref{rmk:singsmall}, singularities arise from `subnetworks'. Since  $\mathcal{M}_{\mathbf{d}, \mathbf{k}, \mathbf{s}, r}^\mathbb{C}$ is the affine cone of its projectification, the same description of singularities holds for the non-projective neuromanifold, apart from the base of the cone $\varphi_\mathbf{w} = 0$, which is a non-nodal singularity. We summarize this in the following corollary. 
\begin{cor}\label{cor:singpts}
Assume that $r>1$. A point $\varphi_\mathbf{w} \in \mathcal{M}_{\mathbf{d}, \mathbf{k}, \mathbf{s}, r}^\mathbb{C}$ is singular if, and only if, $\varphi_\mathbf{w} = 0$ or the fiber of $\varphi_\mathbf{w}$ via $\overline{\varphi}$ contains multiple points. 
%In the latter case, the singularity is nodal and 
In particular, singular points correspond to CNNs that can be parametrized with a smaller architecture.
\end{cor}






% Actually, it is possible to describe the locus where  $\overline{\varphi}$ is an isomorphism more precisely. By the celebrated Zariski Main Theorem \parencite[Corollary III.11.4]{hartshorne2013algebraic}, a finite birational map is an isomorphism once the pre-image of the singular locus has been removed. Therefore, the restriction of $\overline{\varphi}$ 
% \begin{equation}\label{eq:regiso}
% \overline{\varphi}^{-1}\left((\mathbb{P}\mathcal{M}_{\mathbf{d}, \mathbf{k}, \mathbf{s}, r}^\mathbb{C})_{\textnormal{reg}}\right)\rightarrow (\mathbb{P}\mathcal{M}_{\mathbf{d}, \mathbf{k}, \mathbf{s}, r}^\mathbb{C})_{\textnormal{reg}}    
%  \end{equation}
% is an isomorphism. The same holds in the non-projective context, i.e. for $\varphi$. This will be useful in the following section, where we will discuss the Euclidean distance degree of the neurovariety. 

\begin{rmk}\label{rem:realcase}
As anticipated, we remark that most of the results from this section hold over $\mathbb{R}$ as well. This is the case for Corollary \ref{cor:basepts} and Proposition \ref{prop:converonese} -- since they are purely algebraic -- as well as for Corollary \ref{cor:finiteparam} -- since the dimension is independent of the coefficients -- and for Theorem \ref{thm:reg} and Theorem \ref{thm:mainbir} -- since birationality and rank of differential are invariant with respect to restriction of coefficients. However, the proof of Proposition \ref{prop:zarclosed} does not hold over $\mathbb{R}$, since it leverages on properties requiring algebraic closedness. Yet, in the following, we show that the main results of this section imply that the neuromanifold is still closed in the Zariski topology over $\mathbb{R}$.  
\end{rmk}
\begin{prop}\label{prop:zarboundary}
Assume that $r>1$. Then both $\mathbb{P}\mathcal{M}_{\mathbf{d}, \mathbf{k}, \mathbf{s}, r}$ and $\mathcal{M}_{\mathbf{d}, \mathbf{k}, \mathbf{s}, r}$ are closed in the (real) Zariski topology of their respective ambient spaces. 
\end{prop}
\begin{proof}
See Section \ref{app:zarboundary} in the appendix.
\end{proof}
In particular, the (projective) neuromanifold is a real algebraic variety. This does not hold for linear CNNs ($r=1$), in which case the closure in the Zariski topology adds points to the neuromanifold \parencite{kohn2023function}. 
% In particular, $(\overline{\mathcal{M}}_{\mathbf{d}, \mathbf{k}, \mathbf{s}, r})_{\textnormal{reg}} = (\mathcal{M}_{\mathbf{d}, \mathbf{k}, \mathbf{s}, r})_{\textnormal{reg}}$, which will be useful in the following section.
% \begin{defn}
%     Let \(\mathcal{M} \subset \mathbb{R}^n\) be a semi-algebraic set. The relative boundary of \(\mathcal{M}\) in its Zariski closure \(\overline{\mathcal{M}}\), denoted by \(\partial \mathcal{M}\), is defined as
%     \[
%     \partial \mathcal{M} = \overline{\overline{\mathcal{M}} \setminus \mathcal{M}}^{\operatorname{Euc}} \cap \mathcal{M},
%     \]
%     where the last closure is taken in the Euclidean topology.
% \end{defn}

% \begin{lem}
%     Let \( f: \mathcal{X} \to \mathcal{Y} \) be a birational projective morphism between real irreducible varieties \( \mathcal{X} \) and \( \mathcal{Y} \). Then all the points in \( \mathcal{Y} \setminus f(\mathcal{X}) \) and in the relative boundary \( \partial f(\mathcal{X})= \overline{(\mathcal{Y}\setminus f(\mathcal{X}))}^{\operatorname{Euc}} \cap f(\mathcal{X})\) belong to the singular locus of \( \mathcal{Y} \).
% \end{lem}
% \begin{proof}
% Since \( f \) is projective, \( f(\mathcal{X}) \) is closed in the Euclidean topology. Also, \( f \) being birational implies there exist open sets \( \mathcal{U}_{\mathcal{X}} \subset \mathcal{X} \) and \( \mathcal{U}_{\mathcal{Y}} \subset \mathcal{Y} \) such that \( f|_{\mathcal{U}_{\mathcal{X}}} \) is an isomorphism. Since \(\mathcal{U}_{\mathcal{Y}}\) is dense in \( \mathcal{Y} \), we conclude that \( \mathcal{Y} \setminus \mathcal{U}_{\mathcal{Y}} \) has codimension at least 1.
%     Now, let \( y \in \mathcal{Y} \setminus f(\mathcal{X}) \). Since \( f(\mathcal{X}) \) is closed in the Euclidean topology, there exists a Euclidean neighborhood \( \mathcal{N}_y(\mathcal{Y}) \subset \mathcal{Y} \) containing $y$ such that \( \mathcal{N}_y(\mathcal{Y}) \cap f(\mathcal{X}) = \emptyset \). Since the dimension of \( \mathcal{N}_y(\mathcal{Y}) \) is less than \( \dim(\mathcal{Y}) \), we conclude that \( y \) is a singular point of \( \mathcal{Y} \).
%     Finally, the singular locus of a variety is closed in the Zariski topology, and hence it is also closed in the Euclidean topology. Thus, all the points in the relative boundary \( \partial f(\mathcal{X}) \) belong to the singular locus of \( \mathcal{Y} \).
% \end{proof}
\subsection{Optimization}\label{sec:optim}
We now discuss aspects of optimization of a polynomial CNN for a regression task. In particular, by leveraging on the theory of the Euclidean distance degree introduced in Section \ref{sec:eucdd}, we compute an upper bound for the number of (complex) critical points over the neuromanifold for the regression objective. 

We start by introducing the regression objective. Consider a dataset, i.e., a finite subset $\mathcal{D} \subset \mathbb{R}^{d_0} \times \mathbb{R}^{d_L}$ representing input-output pairs. The square-error loss $\mathcal{L}_{\mathcal{D}}\colon \mathcal{M}_{\mathbf{d}, \mathbf{k}, \mathbf{s}, r} \rightarrow \mathbb{R}_{\geq 0}$ is given by: 
\begin{equation}
    \label{eq:loss}
    \mathcal{L}_\mathcal{D}(\varphi_\mathbf{w}) = \sum_{(x,y) \in \mathcal{D}} \| \varphi_{\mathbf{w}}(x) - y \|^2.
\end{equation}
Training a CNN on the dataset $\mathcal{D}$ amounts to minimizing $\mathcal{L}_\mathcal{D}$ over $\mathcal{M}_{\mathbf{d}, \mathbf{k}, \mathbf{s}, r}$. Intuitively, $\varphi_{\mathbf{w}}$ is optimized to interpolate $\mathcal{D}$, which is reminiscent of a closest-point problem over the neuromanifold, but for a discrete set $\mathcal{D}$ instead of an anchor function. Indeed, let $X$ and $Y$ be the matrices whose columns are, respectively, $\nu_{r^{L-1}, d_0 -1}(x) \in \symm{r^{L-1}}{d_0}$ and $y \in \mathbb{R}^{d_L}$ for $(x,y) \in \mathcal{D}$ (in some fixed order), where $\nu$ is the Veronese embedding. It follows that $X$ has full rank for a sufficiently large $|\mathcal{D}|$. Consider the positive definite symmetric matrix $A_\mathcal{D} = XX^\top \otimes I_{d_L}$. By \cite[Section 6]{kohn2022geometry}, we have
\begin{equation}
\label{eq:loss_dist}
 \mathcal{L}_\mathcal{D} (\varphi_\mathbf{w}) =  \textnormal{d}_{A_\mathcal{D}}(\varphi_\mathbf{w}, v_\mathcal{D})^2 + \textnormal{const},
\end{equation}
where $v_\mathcal{D} = YX^\top A_\mathcal{D}^{-1}$. The latter is generic for generic $\mathcal{D}$, since $Y$ can vary independently of $X$. 


% In the following, we show that for large datasets, this problem actually coincides with a closest-point problem. 
% \begin{thm}\label{thm:eddred}
% There exists a linear subspace $V \subset \symm{r^{L-1}}{d_0}^{d_L}$ containing the neuromanifold with the following property. For $|\mathcal{D}| \gg 0$, there exists a positive-definite quadratic form $A_\mathcal{D}$ over $V$ -- inducing a distance $\textnormal{d}_{A_\mathcal{D}}$ -- and $u_\mathcal{D} \in  V$ such that $\mathcal{L}_\mathcal{D}$ and $\textnormal{d}_{A_\mathcal{D}}(\bullet, u_\mathcal{D})^2$ coincide up to an additive constant. Moreover, $u_\mathcal{D}$ is generic for generic $\mathcal{D}$ of fixed cardinality. 
% \end{thm}
% \begin{proof}   
% Let $X$ and $Y$ be the matrices whose columns are, respectively, $\nu_{r^{L-1}, d_0 -1}(x) \in \symm{r^{L-1}}{d_0}$ and $y \in \mathbb{R}^{d_L}$ for $(x,y) \in \mathcal{D}$ (in some fixed order), where $\nu$ is the Veronese embedding. It follows that $X$ has full rank for a sufficiently large $|\mathcal{D}|$. Consider the positive definite symmetric matrix $A_\mathcal{D} = XX^\top \otimes I_{d_L}$. By \cite[Section 6]{kohn2022geometry}, we have
% \begin{equation}
% \label{eq:loss_dist}
%  \mathcal{L}_\mathcal{D} (\varphi_\mathbf{w}) =  \textnormal{d}_{A_\mathcal{D}}(\varphi_\mathbf{w}, v_\mathcal{D})^2 + \textnormal{const},
% \end{equation}
% where $v_\mathcal{D} = YX^\top A_\mathcal{D}^{-1}$. Note that for generic $\mathcal{D}$, the quadratic form induced by $A_\mathcal{D}$ is not generic for $d_L > 1$.  However, by Proposition \ref{prop:converonese}, we can write 
% \begin{equation}\label{eq:genericconv}
% \varphi_{\mathbf{w}}(x) = \tilde{w}~\star_{\tilde{s}}\tilde{x},
% \end{equation}
% where $\tilde{x}$ consists of (possibly repeated) monomials in $x$. Let $V$ the space of linear maps $\symm{r^{L-1}}{d_0} \to \mathbb{R}^{d_L}$ that can be written in the form of Equation \ref{eq:genericconv} for some filter $\tilde{w}$. Such linear maps are determined by one of their output coordinates, meaning that the projections $V \rightarrow \symm{r^{L-1}}{d_0}$ are injective. Now, $A_\mathcal{D}$ contains as a tensor factor the quadratic form $XX^\top$, which is generic over $\symm{r^{L-1}}{d_0}$. It follows that $A_\mathcal{D}$ is generic when restricted to $V$. Lastly, let $u_\mathcal{D}$ be the orthogonal projection of $v_\mathcal{D}$ to $V$ w.r.t. the scalar product induced by $A_\mathcal{D}$. Replacing $v_\mathcal{D}$ by $u_\mathcal{D}$ in Equation \ref{eq:genericconv} introduces an additive constant, which concludes the proof. 
% \end{proof}

% \begin{rmk}
% Similarly to Remark \ref{rmk:basepts}, we point out that the above result is characteristic of CNNs, and does not hold for other architectures. From the proof of Theorem \ref{thm:eddred} it is evident that both $A_\mathcal{D}$ and $u_\mathcal{D}$ can be defined for arbitrary polynomials in $\symm{r^{L-1}}{d_0}^{d_L}$. However, in general, $A_\mathcal{D}$ is not generic, since it contains a tensor factor equal to the identity on the output space $\mathbb{R}^{d_L}$. For CNNs, this is circumvented by exploiting the fact that the first output entry determines the other ones, which follows from the equivariance properties of convolutions. In general, this does not hold. For example, for fully-connected polynomial networks, $\mathcal{A}_\mathcal{D}$ is generic only if $d_L = 1$ -- i.e., for scalar-valued networks -- and it is non-generic otherwise \parencite{kubjas2024geometry}. In the latter case, the theory of the generic Euclidean distance degree is, unfortunately, inapplicable.
% \end{rmk}

% Theorem \ref{thm:eddred} draws a connection between the training of polynomial CNNs and (generic) distance minimization. 

Motivated by this, we compute the generic Euclidean distance degree of the neuromanifold. 
\begin{prop}\label{prop:eddneurovar}
    Assume that $r > 1$ and let $\overline{k} = \textnormal{dim}(\mathcal{M}_{\mathbf{d}, \mathbf{k}, \mathbf{s}, r} )-1 = |\mathbf{k}| - L$. The generic Euclidean distance degree of the neuromanifold is given by: 
    \begin{align*}
         \edd\left(\mathcal{M}_{\mathbf{d}, \mathbf{k}, \mathbf{s}, r} \right) = \sum_{0 \leq i \leq\overline{k} } (-1)^i \left( 2^{\overline{k}  + 1 - i} - 1 \right)(\overline{k}  - i)!  \\ \sum_{\substack{|\mathbf{\alpha}|=i \\ \forall j \  \alpha_j < k_j}}  
         \prod_{0 \leq j < L} \frac{\binom{k_j}{\alpha_j}}{(k_j - \alpha_j - 1)!} \ r^{(L - j -1)(k_j - \alpha_j - 1)},
    \end{align*}
\end{prop}
\begin{proof}
See Section \ref{app:eddneurovar} in the appendix. 
\end{proof}
The above formula is challenging to analyze; numerical values are provided in Table \ref{tab:your_label}. Since it is known that for a projective variety $X$, $\textnormal{gED}(X) \geq \textnormal{deg}(X)$ \parencite[Theorem 5.4]{draisma2016euclidean}, a lower bound follows from  Corollary \ref{cor:finiteparam}. In particular, the growth is super-exponential w.r.t. the network's depth $L$. From the discussion in Section \ref{sec:eucdd},  
it follows that our formula upper-bounds the number of critical points over $(\mathcal{M}_{\mathbf{d}, \mathbf{k}, \mathbf{s}, r}^\mathbb{C})_\textnormal{reg}$ of the (complexified) loss function $\mathcal{L}_ \mathcal{D}$ for large generic datasets $\mathcal{D}$. Since real critical points over a variety remain such after complexification, the upper bound also holds for the critical points of the loss function over $(\mathcal{M}_{\mathbf{d}, \mathbf{k}, \mathbf{s}, r})_{\textnormal{reg}}$. 

A subtle point is that, in practice, optimization is performed via gradient descent over the parameter space, and not over the neuromanifold. The gradient flow of the loss -- and in particular the number of critical points -- could in principle differ due to the parametrization. However, this is not the case for polynomial CNNs: since the projectified parametrization is regular (Theorem \ref{thm:reg}), $\mathbf{w}$ is critical for $\mathcal{L}_{\mathcal{D}} \circ \varphi$ if, and only if, $\varphi_\mathbf{w}$ is critical for $\mathcal{L}_{\mathcal{D}}$. Therefore, our formula equivalently upper-bounds the number of critical points of $\mathcal{L}_\mathcal{D} \circ \varphi$ over $\varphi^{-1}((\mathcal{M}_{\mathbf{d}, \mathbf{k}, \mathbf{s}, r}^\mathbb{C})_{\textnormal{reg}}) \subseteq \mathbb{C}^{|\mathbf{k}|}$, up to the fibers of the parametrization, i.e., up to rescaling each filter. 

Lastly, note that so far we have removed the singular points from the neuromanifold in the above constructions. Therefore, we now discuss the role of singular points of the neuromanifold in the optimization. The following result shows that, excluding the trivial vanishing CNN, the (parameters of) such singular points are not critical for a generic distance minimization problem. 
\iffalse\vs{I think this is what we want to prove: Conjecture \ref{conj:reg} implies $\mathbf{w}$ with non-zero filters is critical for $\mathcal{L}_{\mathcal{D}}$ if and only if $\varphi_{\mathbf{w}}$ is critical for $ \varphi_{\mathbf{w}} \mapsto ||\varphi_{\mathbf{w}}(x)-y||^2$, where $\varphi_{\mathbf{w}} \in (\mathcal{M}_{\mathbf{d}, \mathbf{k}, \mathbf{s}, r})_\textnormal{reg}$. In particular, the number of all complex critical points in optimization of $\mathcal{L}_{\mathcal{D}}$ is exactly $\edd\left( \overline{\mathcal{M}}_{\mathbf{d}, \mathbf{k}, \mathbf{s}, r} \right)$ in Proposition \ref{prop:eddneurovar}.
Proof: By \eqref{eq:loss_dist}, minimizing $\mathcal{L}_{\mathcal{D}}$ is equivalent to minimizing $\mu_{\mathcal{D}}$, where $\mu_{\mathcal{D}}(\mathbf{w}) =\textnormal{d}_{A_\mathcal{D}}(\varphi_\mathbf{w}, v_\mathcal{D})^2$. Note that $\mathbf{w}$ is a critical point of $\mu_{\mathcal{D}}$ if and only if
$\varphi_\mathbf{w}-v_\mathcal{D}$ is in the normal space $\operatorname{im}(d_{\mathbf{w}}\mu)^{\perp_{A_{\mathcal{D}}}}$, where $\perp_{A_{\mathcal{D}}}$ is orthogonal complement with respect to the inner product $\langle \cdot,\cdot \rangle_{A_{\mathcal{D}}}$. Now, let $\Theta$ be the set of all tuple of filters with non-zero filters. Then $\Theta$ has dimension at most $\textnormal{dim}(\mathcal{M}_{\mathbf{d}, \mathbf{k}, \mathbf{s}, r})-1+(L-1) = |\mathbf{k}| -1$ which by \cite[Lemma 7.5]{kohn2023function}, we deduce that no $\mathbf{w} \in \Theta$ is critical for $\mu_{\mathcal{D}}$.  
}\fi

\begin{prop}\label{prop:nosing}
Assume that $r>1$. Let $V \subseteq \symm{r^{L-1}}{d_0}^{d_L}$ be a linear subspace containing the neuromanifold, and $A$ be a positive-definite quadratic form on $V$. For a generic element $u \in V$, if $\mathbf{w}  \in \mathbb{R}^{|\mathbf{k}|}$ is critical for $\textnormal{d}_A(\varphi_{\bullet}, u)^2$, then either $\varphi_\mathbf{w} = 0$ or $\varphi_\mathbf{w} \in (\mathcal{M}_{\mathbf{d}, \mathbf{k}, \mathbf{s}, r})_\textnormal{reg}$.
\end{prop}
\begin{proof}
See Section \ref{app:nosing} in the appendix. 
\end{proof}
Therefore, if we exclude vanishing filters, we obtain an upper bound on the number of critical points in parameter space, up to rescaling each filter, of $\mathcal{L}_{\mathcal{D}} \circ \varphi$ for large generic $\mathcal{D}$. 

\section{CONCLUSIONS AND FUTURE WORK}
In this work, we have studied aspects of the geometry and optimization of polynomial CNNs. In particular, we have proven that the projectified parametrization is regular and finite birational, derived the dimension of the neuromanifold, and related the latter to the Segre--Veronese variety. Moreover, we have rephrased the optimization of the regression loss as a distance minimization problem, and leveraged on the theory of the Euclidean distance degree to upper-bound the number of (complex) critical points for a generic large dataset. 

Our tools involve general yet powerful results from algebraic geometry. Therefore, similar arguments might be applicable to other neural architectures. For example, \emph{graph neural networks} \parencite{kipf2016semi, bronstein2021geometric} are nowadays popular in a variety of domains, and are closely related to CNNs. Therefore, exploring applications of algebraic geometry to the neuromanifold of such networks represents an interesting line for future investigation. 

A limitation of the Euclidean distance degree is that it quantifies the number of critical points, without distinguishing local maxima, minima, or saddle points. Since local minima are the stable equilibria of gradient descent to which the latter converges, counting them represents a fundamental challenge from the perspective of optimization and learning. To this end, tools from \emph{Morse theory} \parencite{milnor1963morse} -- an approach relating critical points with prescribed index to topological invariants -- might be applicable to neuromanifolds, providing insights into local minima of the loss function. This constitutes another interesting direction to investigate. 


From a broader perspective, a fundamental challenge lies in extending the study of neuromanifolds beyond the algebraic setting, i.e., for activation functions that are not polynomial (e.g., sigmoid and ReLU). To this end, a promising approach is via polynomial approximation. By the Weierstrass Approximation Theorem, any continuous function can be (locally) approximated by polynomials with arbitrary precision. Therefore, tools from algebraic geometry can be potentially applied beyond the algebraic realm, at least approximately. For CNNs, our work is the first to tackle the polynomial case; extending the results to non-polynomial CNNs (e.g. ReLU CNNs) via approximation is a natural and interesting next step. 

\section*{Acknowledgements}
We thank Rainer Sinn for helpful discussions on real algebraic geometry and the recovery of polynomials from partial terms.
This work was partially supported by the Wallenberg AI, Autonomous Systems and Software Program (WASP) funded by the Knut and Alice Wallenberg Foundation. 


\printbibliography 

% \section*{Checklist}
%  \begin{enumerate}
% \item For all models and algorithms presented, check if you include:
% \begin{enumerate}
% \item A clear description of the mathematical setting, assumptions, algorithm, and/or model. \\
% Yes. The paper is completely formal, and all the necessary mathematical background is introduced in Section \ref{sec:back}. 
% \item An analysis of the properties and complexity (time, space, sample size) of any algorithm. \\
% Not Applicable. The paper does not include any algorithm.
% \item (Optional) Anonymized source code, with specification of all dependencies, including external libraries. \\
% Not Applicable. The paper does not involve any implementation. 
%  \end{enumerate}


%  \item For any theoretical claim, check if you include:
%  \begin{enumerate}
%    \item Statements of the full set of assumptions of all theoretical results. \\
%    Yes. Every novel mathematical result is stated formally -- see Section \ref{sec:convneu}. 
%    \item Complete proofs of all theoretical results. \\
%     Yes. Every novel mathematical result is proven formally -- see Section \ref{sec:convneu}. 
%    \item Clear explanations of any assumptions. \\
%     Yes. In Section \ref{sec:convneu}, formal statements are followed by intuitive explanations.  
%  \end{enumerate}


%  \item For all figures and tables that present empirical results, check if you include:
%  \begin{enumerate}
%    \item The code, data, and instructions needed to reproduce the main experimental results (either in the supplemental material or as a URL). \\
%    Not Applicable. The paper does not include any experimental result. 
%    \item All the training details (e.g., data splits, hyperparameters, how they were chosen). \\
%       Not Applicable. The paper does not include any experimental result. 
%      \item A clear definition of the specific measure or statistics and error bars (e.g., with respect to the random seed after running experiments multiple times). \\
%     Not Applicable. The paper does not include any experimental result. 
%      \item A description of the computing infrastructure used. (e.g., type of GPUs, internal cluster, or cloud provider). \\
%       Not Applicable. The paper does not include any experimental result. 
%  \end{enumerate}

%  \item If you are using existing assets (e.g., code, data, models) or curating/releasing new assets, check if you include:
%  \begin{enumerate}
%    \item Citations of the creator If your work uses existing assets. \\
%     Not Applicable. The paper does not include any existing asset.
%    \item The license information of the assets, if applicable. \\
%     Not Applicable. The paper does not include any existing asset.
%    \item New assets either in the supplemental material or as a URL, if applicable. \\
%     Not Applicable. The paper does not include any new asset.
%    \item Discussion of sensible content if applicable, e.g., personally identifiable information or offensive content. \\
%     Not Applicable. The paper does not include any asset.
%  \end{enumerate}

%  \item If you used crowdsourcing or conducted research with human subjects, check if you include:
%  \begin{enumerate}
%    \item The full text of instructions given to participants and screenshots. \\
%     Not Applicable. The paper does not include any crowdsourcing or experiments with human subjects.
%    \item Descriptions of potential participant risks, with links to Institutional Review Board (IRB) approvals if applicable.  \\
%     Not Applicable. The paper does not include any crowdsourcing or experiments with human subjects.
%    \item The estimated hourly wage paid to participants and the total amount spent on participant compensation.  \\
%     Not Applicable. The paper does not include any crowdsourcing or experiments with human subjects.
%  \end{enumerate}

%  \end{enumerate}

\newpage 
\appendix 

\onecolumn 
\section{ON THE DIFFERENTIAL OF THE PARAMETRIZATION}\label{sec:proofdet}
In this section, we prove technical results on the differential of the parametrization of polynomial CNNs. The regularity of the projectified parametrization follows immediately -- see Theorem \ref{thm:reg}. We adhere to the convention from Section \ref{sec:convneu} and work with complex scalars. As discussed in Remark \ref{rem:realcase}, the results hold, a posteriori, over $\mathbb{R}$ as well. Moreover, we denote by $\textnormal{J}_\mathbf{w} \varphi$ the differential of $\varphi$ at $\mathbf{w}$. Since both the domain and co-domain of $\varphi$ are vector spaces, the differential can be seen as a linear map $\textnormal{J}_\mathbf{w} \varphi \colon \mathbb{C}^{|\mathbf{k}|} \rightarrow \symm{r^{L-1}}{d_0}^{d_L}$. Moreover, for $L>1$, we denote $\mathbf{w}' = (w_{0}, \ldots, w_{L-2})$. Then $\varphi_{\mathbf{w}} = w_{L-1} \conv{s_{L-1}} \sigma_r(\varphi_{\mathbf{w}'})$, and by applying the Leibniz derivation rule we see that $\textnormal{J}_\mathbf{w} \varphi$ sends a tangent vector $\dot{\mathbf{w}} = (\dot{w}_{0}, \ldots, \dot{w}_{L-1}) \in \mathbb{C}^{|\mathbf{k}|}$ to:
\begin{equation}\label{eq:diffhadam}
\dot{w}_{L-1} \conv{s_{L-1}} \sigma_r\left(  \varphi_{\mathbf{w}'} \right) + r w_{L-1} \conv{s_{L-1}} \sigma_{r-1}\left(  \varphi_{\mathbf{w}'} \right) \odot \left( \textnormal{J}_\mathbf{w'} \varphi \right)(\dot{\mathbf{w}}'),
\end{equation}
where $\odot$ denotes the Hadamard product, i.e., the point-wise product between polynomial functions. 

\begin{lem}\label{lemm:homogendiff}
Given $\lambda_0, \ldots, \lambda_{L-1} \in \mathbb{C}$, the differential $\textnormal{J}_\mathbf{w} \varphi$ sends $(\lambda_{0}w_{0}, \ldots, \lambda_{L-1} w_{L-1})$ to:
\begin{equation}
    (\lambda_{L-1} + r \lambda_{L-2} + \cdots + r^{L-1}\lambda_0) \ \varphi_\mathbf{w}. 
\end{equation}
\end{lem}
\begin{proof}
This follows immediately by substituting $\dot{w}_i = \lambda_i w_i$ for all $i$ in Equation \ref{eq:diffhadam} and by reasoning inductively on the number of layers $L$. 
\end{proof}

\begin{prop}\label{prop:kerdiff}
Suppose that $w_i \not = 0 $ for all $i$. Then the kernel of $\textnormal{J}_\mathbf{w} \varphi$ is generated, as a linear subspace of $\mathbb{C}^{|\mathbf{k}|}$, by $(0,\ldots,   0, w_{L-2}, -rw_{L-1})$, $(0, \ldots, 0, w_{L-3}, -rw_{L-2},0)$, $\ldots$, $(w_0, -rw_1, 0, \ldots, 0)$. 
\end{prop}
\begin{proof}
We proceed by induction over the number of layers $L$. For $L=1$,  $\varphi_\mathbf{w}(x) = w_0 \conv{s_0} x$, which is linear in $w_0$ and does not vanish for all $x$ since $w_0 \not = 0$. The differential is therefore injective, and its kernel is trivial. 

Suppose now that $L > 1$. In order to compute the kernel $\textnormal{J}_\mathbf{w} \varphi$, we need to solve
\begin{equation}\label{eq:kerndiff}
   \dot{w}_{L-1} \conv{s_{L-1}} \sigma_r\left(  \varphi_{\mathbf{w}'} \right) + r w_{L-1} \conv{s_{L-1}} \sigma_{r-1}\left(  \varphi_{\mathbf{w}'} \right) \odot \left( \textnormal{J}_\mathbf{w'} \varphi \right)(\dot{\mathbf{w}}') = 0. 
\end{equation}
If $\dot{\mathbf{w}}'$ belongs to the kernel of $\textnormal{J}_\mathbf{w'} \varphi$, then since $\varphi_{\mathbf{w}'} \not = 0$ by hypothesis (see Corollary \ref{cor:basepts}), we must have $\dot{w}_{L-1} = 0$. It follows from the inductive hypothesis that $\dot{w}$ is a linear combination of the tuples of filters in the statement. We therefore assume that $\left( \textnormal{J}_\mathbf{w'} \varphi\right)(\dot{\mathbf{w}}') \not = 0$. Moreover, without loss of generality, we assume that  $w_{L-1}[0] \not = 0$. 

We will now decompose $\varphi_\mathbf{w}$ as a sum of CNNs with smaller filter size on the last layer. To this end, consider the tuple of filters:
\begin{equation}
\mathbf{w}^+=(w_{0}, \ldots, w_{L-2}, (w_{L-1}[0])), \hspace{4em} \mathbf{w}^-=(w_{0}, \ldots, w_{L-2}, (0, w_{L-1}[1], \ldots, w_{L-1}[k_{L-1} -1])).
\end{equation}
Then $\varphi_\mathbf{w} = \varphi_{\mathbf{w}^+} + \varphi_{\mathbf{w}^-}$. Now, let $i$ be the minimal index such that $x[i]$ appears in $\varphi_{\mathbf{w}'}[0]$, seen as a scalar-valued polynomial in the input variable $x$. Recall the basic equivariance property of CNNs, meaning that shifting indices in the output variable is equivalent to shifting them in the input one (assuming the latter shift accounts for the strides). This implies that $x[i]$ does not appear in neither $\left( \textnormal{J}_\mathbf{w^-} \varphi \right)(\dot{\mathbf{w}}^-)$ nor $\varphi_{\mathbf{w}^-}$. Since $0 = \textnormal{J}_\mathbf{w} \varphi = \textnormal{J}_{\mathbf{w}^+} \varphi + \textnormal{J}_{\mathbf{w}^-} \varphi$, $x[i]$ does not appear in $\textnormal{J}_{\mathbf{w}^+} \varphi$ as well. Write: 
\begin{equation}
\varphi_{\mathbf{w}'}[0] = x[i]f + g, \hspace{7em} \left( \textnormal{J}_\mathbf{w'} \varphi \right)(\dot{\mathbf{w}}')[0] = x[i]a + b,
\end{equation}
where $f,g,a,b$ are polynomials such that $f \not =  0$ and $x[i]$ does not appear in $g$ nor $b$. The (vanishing) terms containing $x[i]$ in $\textnormal{J}_{\mathbf{w}^+} \varphi$ have the form:
\begin{equation}
    0 = \dot{w}_{L-1}[0] \left((x[i]f + g)^r - g^r \right) + rw_{L-1}[0]\left((x[i]f + g)^{r-1}(x[i]a + b) - g^{r-1}b\right).
\end{equation}
By manipulating the above expression, we obtain:
\begin{equation}
(x[i]f + g)^{r-1}\left(\dot{w}_{L-1}[0](x[i]f + g) + rw_{L-1}[0](x[i]a + b) \right) = g^{r-1}( \dot{w}_{L-1}[0]g + r w_{L-1}[0]b).
\end{equation}
The right-hand side of the above equation does not contain $x[i]$. Since $f \not = 0$, it follows that $\dot{w}_{L-1}[0](x[i]f + g) + rw_{L-1}[0](x[i]a + b) = 0$, implying:
\begin{equation}
\varphi_{\mathbf{w}'}[0] = x[i]a + b = \lambda \ (x[i]f + g) = \lambda \left( \textnormal{J}_\mathbf{w'} \varphi \right)(\dot{\mathbf{w}}')[0], \hspace{3em} \lambda := -\frac{\dot{w}_{L-1}[0]}{w_{L-1}[0]}.
\end{equation}
By looking at (the $0$-th entry of) Equation \ref{eq:kerndiff}, we deduce $\dot{w}_{L-1} = -r \lambda w_{L-1}$. By Lemma \ref{lemm:homogendiff}, $\dot{\mathbf{w}}' = (\lambda_{0}w_{0}, \ldots, \lambda_{L-2}w_{L-2})$ for some $\lambda_i \in \mathbb{C} \setminus \{0\}$ such that $\lambda_{L-2} + r \lambda_{r-3} + \cdots + r^{L-2} \lambda_0 = \lambda$. But then
\begin{align}
\dot{\mathbf{w}} &= (\lambda_0 w_0, \ldots,  \lambda_{L-2}w_{L-2}, -r \lambda w_{L-1}) = \\
&= \lambda (0, \ldots, 0, w_{L-2}, -rw_{L-1}) + (\lambda_{0} w_0, \ldots, \lambda_{L-3}w_{L-3}, (\lambda_{L-2} - \lambda)w_{L-2}, 0).
\end{align}
Since $(\lambda_{L-2} - \lambda) + r \lambda_{L-3} + \cdots + r^{L-2}\lambda_0 = 0$, it follows from Lemma \ref{lemm:homogendiff} that the second summand of the above expression belongs to the kernel of $ \textnormal{J}_\mathbf{w'} \varphi$, and the desired result follows from the inductive hypothesis.
\end{proof}



\section{ADDITIONAL PROOFS}\label{sec:remproof}
\subsection{Proof of Lemma \ref{lem:fullrk}}\label{app:fullrk}
\begin{proof}
Since the domain of a convolution has higher dimension than the co-domain, we need to show that the rows $A_0, \ldots, A_{d'-1}$ of the corresponding Toeplitz matrix are linearly independent. By induction on the filter size $k$, we can assume $w[0] \not = 0$. Suppose that $\sum_i a_i A_i = 0$ for some scalars $a_0, \ldots, a_{d'-1}$. Since the first column of the Toeplitz matrix has only one non-vanishing entry, we deduce $a_0 = 0$. But then, since the $s$-th column is $(w[s-1], w[0], 0, \ldots, 0)$, we similarly deduce that $a_1 = 0$, and so on. 
\end{proof}

\subsection{Proof of Proposition \ref{prop:zarclosed}}\label{app:zarclosed}
\begin{proof}
Projective morphisms over algebraically-closed fields are closed \parencite[II, \S 4, Theorem 4.9]{hartshorne2013algebraic} and, in particular, have closed image. Therefore, $\mathbb{P}\mathcal{M}_{\mathbf{d}, \mathbf{k}, \mathbf{s}, r}^\mathbb{C}$ is closed in $\mathbb{P} \symmm{r^{L-1}}{\mathbb{C}^{d_0}}^{d_L}$. Since $\varphi$ is homogeneous, the neuromanifold coincides with the affine cone of its projectification. Since affine cones of varieties are varieties, $\mathcal{M}_{\mathbf{d}, \mathbf{k}, \mathbf{s}, r}^\mathbb{C}$ is closed in $\symmm{r^{L-1}}{\mathbb{C}^{d_0}}^{d_L}$.     
\end{proof}

\subsection{Proof of Proposition \ref{prop:converonese}}\label{app:converonese}
\begin{proof} 
 Newton's multinomial expansion yields: 
  \begin{equation*}\label{eq:newtexp}
      \displaystyle (w \conv{s} x)[i]^r = \sum_{|\mathbf{a}| = r}\binom{r}{\mathbf{a}} \prod_{0 \leq j < k} w[j]^{a_j} \prod_{0 \leq j < k}x[is + j]^{a_j},
\end{equation*}
where $|\mathbf{a}| = a_0 + \cdots + a_{k-1}$ and $\binom{r}{\mathbf{a}} = \frac{r!}{a_0! \cdots a_{k-1}!}$. We construct $\widetilde{x}$ by indexing its coordinates via pairs $(i, \mathbf{a})$, where $1 \leq i < d - k$ and $\mathbf{a}$ is a multi-index with $|\mathbf{a}| = r$, in some fixed order. The corresponding coordinate of $\tilde{x}$ is $\prod_j x[i +j]^{a_j}$. Similarly, $\tilde{w}$ is indexed by $\mathbf{a}$, with corresponding coordinate $\binom{r}{\mathbf{a}}\prod_j w[j]^{a_j}$, i.e., $\tilde{w}$ is a rescaling of the Veronese embedding of $w$ of degree $r$. The claim follows immediately. 
\end{proof}


\subsection{Proof of Theorem \ref{thm:mainbir}}\label{app:mainbir}
% Before proving the result, we remark a general fact on polynomial CNNs. Consider $\varphi_\mathbf{w}(x)[0]$, seen as a (homogeneous) polynomial in the input variable $x$. Then, given $0 \leq i < d_0$, the coefficient of the monomial $x[i] x[0]^{r^{L-1}-1}$ -- which is a (homogeneous) polynomial in $\mathbf{w}$ -- is, up to multiplicative scalar, of the form:
% \begin{equation}\label{eq:monofib}
% w_0[0]^{r^{L-1} - 1} w_1[0]^{r^{L-2}}  w_2[0]^{r^{L-3}} \cdots w_{L-1}[0]  w_0[i]. 
% \end{equation}
% This can be seen by direct computation. In short, by definition of convolution, the variable $x[0]$ appears in the output of each layer only at the $0$-th entry. Therefore, $x[i]$ can appear as a linear factor only accompanied by its corresponding weight at the first layer, i.e., $w_0[i]$. 
% \begin{proof}
% We will prove an equivalent statement on the non-projectified parameterization. Suppose that $\varphi_\mathbf{w} = \varphi_\mathbf{v}$ for some $\mathbf{w}, \mathbf{v}$ such that $w_i, v_i \not = 0$ for all $0 \leq i < L$. We wish to show that the filters in $\mathbf{w}$ are proportional to the ones in $\mathbf{v}$, if $w_i[0],v_i[0] \not = 0$ for all $i$. To this end, we proceed by induction on $L$. For $L = 1$, the desired result follows from the fact that the linear map defined by a convolution determines the corresponding filter. Therefore, assume $L > 1$. By considering the expression from Equation \ref{eq:monofib} -- which, up to scalar, corresponds to the coefficient of a monomial in $\varphi_\mathbf{w}[0]$ -- we conclude that for all $0 \leq i < k_0$:
% \begin{equation}
% \frac{w_0[i]}{v_0[i]} = \frac{v_0[0]^{r^{L-1} - 1} v_1[0]^{r^{L-2}}  v_2[0]^{r^{L-3}} \cdots v_{L-1}[0]}{w_0[0]^{r^{L-1} - 1} w_1[0]^{r^{L-2}}  w_2[0]^{r^{L-3}} \cdots w_{L-1}[0]} := \lambda.
% \end{equation}
% Therefore, $w_0$ is proportional to $v_0$. In particular, $\varphi_{\mathbf{w}'}(\lambda^r \ \sigma_r(v_0 \conv{s_0} x)) = \varphi_{\mathbf{v}'}(\sigma_r(v_0 \conv{s_0} x)) $, where $\mathbf{w}' = (w_1, \ldots, w_{L-1})$ denotes the $(L-1)$-uple of filters where the first one has been removed, and similarly for $\mathbf{v}'$. Since by Lemma \ref{lem:fullrk} the map $x \mapsto \sigma_r(v_0 \conv{s_0} x)$ is surjective, we conclude that $\varphi_{\mathbf{w}'}$ is proportional to $\varphi_{\mathbf{v}'}$, or, equivalently, they coincide up to rescaling one of the filters. By the inductive hypothesis, $w_i$ is proportional to $v_i$ for all $i > 0$, as desired.   
% % We now aim to characterize the special fibers. Given $\mathbf{w}$, we wish to prove that if $w_i[0] = 0$ for some $i$, then there is a tuple of filters $\mathbf{v}$ that is \emph{not} proportional to $\mathbf{w}$ and such that $\varphi_\mathbf{w} = \varphi_{\mathbf{v}}$. Such $\mathbf{v}$ can be constructed as follows. Starting from $\mathbf{v} = \mathbf{w}$, pick an arbitrary $j \not = 0$, and then set $v_j[0] = 0$ if $w_j[0] \not = 0$, and 
% \end{proof}
We first provide a general result on homogeneous polynomials.
\begin{lem}\label{lemm:eulerproj}
Assume that $r > 1$. Let $d,n \in \mathbb{N}$, $0 \leq i < n$, and $p$ be a multivariate homogeneous polynomial of degree $d$ in $n$ variables over a field of characteristic $0$ (e.g., $\mathbb{R}$ or $\mathbb{C}$). Then the terms containing $x[i]$ in $p^r$ determine $p$ up to multiplicative scalar. 
\end{lem}
\begin{proof}
Write $p = \sum_{0 \leq j \leq d}x[i]^jq_j$, where $q_j$ is a homogeneous polynomial of degree $d-j$ not containing $x[i]$. Newton's multinomial expansion yields:
\begin{equation}
p^r = \sum_{|\mathbf{a}| = r}\binom{r}{\mathbf{a}} \prod_{0 \leq j \leq d} x[i]^{j a_j}q_j^{a_j} =  \sum_{|\mathbf{a}| = r}\binom{r}{\mathbf{a}} x[i]^{\sum_{j}j a_j} \prod_{0 \leq j \leq d} q_j^{a_j} .  
\end{equation}
By following $t := \sum_{0 \leq j \leq d}ja_j$ in decreasing order, we can inductively recover $q_j$ up to scalar in decreasing order w.r.t. $j$. More specifically, starting from $t = dr$, the coefficient of $x[i]^t$ is $q_d^{d}$, from which we recover $q_d$. Then, for $t = dr -1$, the coefficient of $x[i]^t$ is $q_d^{r-1}q_{d-1}$, from which we can now recover $q_{d-1}$, and so on.
\end{proof}
We are now ready to prove Theorem \ref{thm:mainbir}.
\begin{proof}
We will prove the following equivalent statement on the non-projectified parametrization: if $\varphi_\mathbf{w}[0] = \varphi_\mathbf{v}[k]$ for some $0 < k < d_L$ and some $\mathbf{w}, \mathbf{v} \in \mathbb{C}^{|\mathbf{k}|}$ such that $w_i, v_i \not = 0$ for all $0 \leq i < L$, then $\mathbf{w}$ and $\mathbf{v}$ are related, up to rescaling, by a shift of their indices. To this end, we proceed by induction on $L$. For $L = 1$, the desired result follows from the fact that the linear map defined by a convolution determines the corresponding filter. Suppose now that $L > 1$ and that $\varphi_\mathbf{w}[0] = \varphi_\mathbf{v}[k]$. The inductive argument is similar in spirit to the one from the proof of Proposition \ref{prop:kerdiff}. Denote  $\mathbf{w}' = (w_{0}, \ldots, w_{L-2})$, and analogously for $\mathbf{v}'$. Then for all $x$: 
\begin{equation}\label{eq:shiftfil}
\varphi_\mathbf{w}(x)[0] = \sum_{0 \leq j < k_{L-1}} w_{L-1}[j]\ (\varphi_{\mathbf{w}'}(x)[j])^r = \varphi_\mathbf{v}(x)[k] = \sum_{0 \leq j < k_{L-1}} v_{L-1}[j]\ (\varphi_{\mathbf{v}'}(x)[s_{L-1}k + j])^r. 
\end{equation}
Let $i$ be the minimal index such that $x[i]$ appears in the above expression, seen as a scalar-valued polynomial in the input variable $x$. Moreover, let $m$ and $n$ be the minimal indices such that $w_{L-1}[m] \not = 0$ and $v_{L-1}[n] \not = 0$, respectively. Note that among all the summands in Equation \ref{eq:shiftfil}, $x[i]$ appears only in $w_{L-1}[m] \ (\varphi_{\mathbf{w}'}(x)[m])^r $ and $v_{L-1}[n] \ (\varphi_{\mathbf{v}'}(x)[s_{L-1}k + n])^r$. Therefore, the terms involving $x[i]$ in these two summands must coincide. By Lemma \ref{lemm:eulerproj}, $\varphi_{\mathbf{w}'}(x)[m] = \lambda \ \varphi_{\mathbf{v}'}(x)[s_{L-1}k + n]$ for some $\lambda \in \mathbb{C} \setminus \{ 0\}$. From the equivariance properties of convolutions, it follows that $\varphi_{\mathbf{w}'}(x)[j] = \lambda \ \varphi_{\mathbf{v}'}(x)[s_{L-1}k + n - m + j]$ for all $0 \leq j < d_{L-2}$. From the inductive hypothesis, for all $0\leq t < L-1$, $w_t$ and $v_t$ are related, up to rescaling, by a shift of their indices.  Moreover, Equation \ref{eq:shiftfil} expands to:
\begin{equation}\label{eq:sumindxs}
\sum_{n \leq j < k_{L-1}} v_{L-1}[j] \ (\varphi_{\mathbf{v}'}(x)[s_{L-1}k + j])^r = \sum_{m \leq j < k_{L-1}} w_{L-1}[j] \lambda^r \  (\varphi_{\mathbf{v}'}(x)[s_{L-1}k + n - m + j])^r.
\end{equation}
Again, in the above expression $x[i]$ appears only in the summands $v_{L-1}[n]\ (\varphi_{\mathbf{v}'}(x)[s_{L-1}k + n])^r $ and $w_{L-1}[m] \lambda^r \  (\varphi_{\mathbf{v}'}(x)[s_{L-1}k + n])^r$. Since these summands must coincide, we conclude that $v_{L-1}[n] = w_{L-1}[m]\lambda^r$. But then the indices of the sums in Equation \ref{eq:sumindxs} start from $n+1$ and $m+1$. By iterating this argument, we conclude that $v_{L-1}[n - m + j] = w_{L-1}[j]\lambda^r$ for all $j$, i.e., the filters of the last layer coincide up to rescaling and shifting the indices, as desired. 

Lastly, note that the birationality statement on $\overline{\varphi}$ follows immediately. Indeed, the above argument shows that the fiber of $\overline{\varphi}$ at $\varphi_{\mathbf{w}}$ is not a singleton only if some filters in $\mathbf{w}$ are padded by zeros on the left or on the right, which is a negligible condition, i.e., it does not hold almost everywhere. The remaining fibers are characterized by shifting indices of filters, and are therefore finite. 
\end{proof} 



\subsection{Proof of Corollary \ref{cor:finiteparam}}\label{app:finiteparam}
\begin{proof}
Birational maps preserve the dimension. Therefore, the dimension of the projective neuromanifold is:
\begin{equation}
\textnormal{dim}(\mathbb{P}\mathcal{M}_{\mathbf{d}, \mathbf{k}, \mathbf{s}, r}^\mathbb{C}) = \dim (\mathbb{P} \mathbb{C}^{k_0} \times \cdots \times \mathbb{P} \mathbb{C}^{k_{L-1}}) = \sum_{0 \leq i < L} (k_i - 1) = |\mathbf{k}| - L.
\end{equation}
  Since $\textnormal{dim}(\mathbb{P}\mathcal{M}_{\mathbf{d}, \mathbf{k}, \mathbf{s}, r}^\mathbb{C}) = \textnormal{dim}(\mathcal{M}_{\mathbf{d}, \mathbf{k}, \mathbf{s}, r}^\mathbb{C}) -1$,  the dimension of the neuromanifold is $|\mathbf{k}| - L + 1$.

A birational linear (projective) map from an algebraic variety whose kernel does not intersect the variety preserves the degree. Since the projective neuromanifold is the image of a Segre--Veronese variety via a birational linear map, we have that $\deg(\mathbb{P}\mathcal{M}_{\mathbf{d}, \mathbf{k}, \mathbf{s}, r}^\mathbb{C}) = \deg(\mathcal{V}_{\mathbf{m}, \mathbf{p}})$. The latter can be computed via  \cite[Proposition 6.11]{kozhasov2023minimal}, as follows:
\begin{equation}
    \label{eq:deg_neuro}
    \deg(\mathcal{V}_{\mathbf{m}, \mathbf{p}}) = (|\mathbf{k}|-L)!\prod_{0 \leq j < L} \frac{r^{(L - j - 1)(k_j - 1)}}{(k_j - 1)!}.
\end{equation}
Finally, since the degree of a projective variety coincides with the one of its affine cone, we have $\deg(\mathcal{M}_{\mathbf{d}, \mathbf{k}, \mathbf{s}, r}^\mathbb{C}) = \deg(\mathbb{P}\mathcal{M}_{\mathbf{d}, \mathbf{k}, \mathbf{s}, r}^\mathbb{C})$. 
\end{proof}


\subsection{Proof of Proposition \ref{prop:zarboundary}}\label{app:zarboundary}
\begin{proof}
This follows from the regularity of the parametrization. Specifically, for an image of a real algebraic map, the relative boundary of the image inside its Zariski closure is contained in the branch locus of the complexification of the map. Since the complexified projectified parametrization $\overline{\varphi}$ is regular (Theorem \ref{thm:reg}), it is unramified, i.e., the branching locus is empty. Therefore, $\mathbb{P}\mathcal{M}_{\mathbf{d}, \mathbf{k}, \mathbf{s}, r}$ is closed in the Zariski topology, and the same holds for its affine cone $\mathcal{M}_{\mathbf{d}, \mathbf{k}, \mathbf{s}, r}$.
 \end{proof}


\subsection{Proof of Proposition \ref{prop:eddneurovar}}\label{app:eddneurovar}
\begin{proof}
From Section \ref{sec:geome}, we know that $\overline{\varphi}$ factors into the Segre--Veronese embedding $\nu_{\mathbf{m}, \mathbf{p}}$, with $\mathbf{m} = (r^{L-1}, \ldots, r, 1)$ and $\mathbf{p} = (k_0 - 1, \ldots, k_{L-1} - 1)$, followed by a linear map. By \cite[Corollary 6.1]{draisma2016euclidean}, linear projections of varieties of co-dimension $\geq 2$ do not alter the generic Euclidean distance degree. Therefore, $\edd(\mathbb{P}\mathcal{M}_{\mathbf{d}, \mathbf{k}, \mathbf{s}, r}) = \edd(\mathcal{V}_{\mathbf{m}, \mathbf{p}})$. The generic Euclidean distance degree of the Segre--Veronese variety is given by Theorem \ref{th:segverdeg}, from which our formula follows via elementary algebraic manipulations. Note that, by definition, the Euclidean distance degree of a projective variety coincides with the one of its affine cone. Since $\varphi$ is homogeneous, the affine cone of $\mathbb{P}\mathcal{M}_{\mathbf{d}, \mathbf{k}, \mathbf{s}, r}$ is $\overline{\mathcal{M}}_{\mathbf{d}, \mathbf{k}, \mathbf{s}, r}$, concluding the proof.     
\end{proof}

Below, we provide a table with numerical values of the generic Euclidean distance degree. 
\begin{table}[!h]
    \centering
    \renewcommand{\arraystretch}{1.5}
    \setlength{\tabcolsep}{12pt} % Adjust column width
    \begin{tabular}{c|ccccc}
        \diagbox[width=6em, height=2em]{\footnotesize $r$}{\footnotesize $k$} & \footnotesize $2$ & \footnotesize $3$ & \footnotesize $4$ & \footnotesize $5$ & \footnotesize $6$ \\
        \hline
        \footnotesize $1$ & 6 & 39 & 284 & 2205 & 17730 \\
        \footnotesize $2$ & 14 & 219 & 3772 & 68405 & 1277898 \\
        \footnotesize $3$ & 22 & 543 & 14684 & 417005 & 12186066 \\
        \footnotesize $4$ & 30 & 1011 & 37244 & 1439205 & 57202074 \\
        \footnotesize $5$ & 38 & 1623 & 75676 & 3699005 & 185917794 \\
        \footnotesize $6$ & 46 & 2379 & 134204 & 7933205 & 482134890 \\
    \end{tabular}
    \caption{The generic Euclidean distance degree of neuromanifolds with $L=2$ and $k:=k_1=k_2$.}
    \label{tab:your_label}
\end{table}

\subsection{Proof of Proposition \ref{prop:nosing}}\label{app:nosing}
\begin{proof}
Note that $\mathbf{w}$ is critical for $\textnormal{d}_A(\varphi_{\bullet}, u)^2$ if, and only if, $\varphi_{\mathbf{w}} - u$ is perpendicular, according to the scalar product induced by $A$, to the image of the differential of $\varphi$ at $\mathbf{w}$. In other words, $u$ must belong to the relative normal bundle of $\varphi$. By Proposition \ref{cor:finiteparam} and Theorem \ref{thm:reg}, such bundle consists, at a given $\varphi_\mathbf{w}$, of a finite number of affine subspaces of co-dimension $|\mathbf{k}| - L + 1 = \textnormal{dim}((\mathcal{M}_{\mathbf{d}, \mathbf{k}, \mathbf{s}, r}))$. If we constrain $\varphi_\mathbf{w}$ to be singular, then the restricted bundle has co-dimension $\textnormal{dim}(\mathcal{M}_{\mathbf{d}, \mathbf{k}, \mathbf{s}, r}) - \textnormal{dim}((\mathcal{M}_{\mathbf{d}, \mathbf{k}, \mathbf{s}, r})_\textnormal{sing}) > 0$. Belonging to it is, therefore, a negligible condition in $u \in V$, which concludes the proof. 
\end{proof}


\section{A Toy Example}\label{sec:toyexmp}
Here, we provide a toy example of a polynomial CNN, where all the quantities from our results can be easily computed. We consider a network with filter sizes $k_0 = k_1 = 2$, stride $s_0 = 1$, and activation function $\sigma_2(x) = x^2$. The function computed by the network is the homogeneous quadratic polynomial:
\begin{equation}
\varphi_\mathbf{w}(x)= a^2 c x_0^2 + 2abc x_0 x_1 + (b^2 c + a^2 d)x_1^2 + 2abd x_1 x_2 + b^2 d x_2^2,
\end{equation} 
where $w_0 = (a,b), w_1 = (c,d)$. The neuromanifold has dimension $3$, degree $4$, and Euclidean distance degree $14$. The singular points, as described in Theorem 4.6, arise when $a=d=0$ or $c=b=0$. These singular points correspond to the self-intersection of the neuromanifold, which in this example is a one-dimensional line.


%From here, we distinguish two cases based on filter size of the last layer:

% \textit{Case 1:} $k_{L-1} = 1$. In this case, the convolutions in Equation \ref{eq:kerndiff} reduce to multiplication by the scalars $\dot{w}_{L-1}$ and $w_{L-1}$, respectively. Therefore, we have that  $\left( \textnormal{J}_\mathbf{w'} \varphi \right)(\dot{\mathbf{w}}') = \lambda \varphi_{\mathbf{w}'}$ and $\dot{w}_{L-1} = -r \lambda w_{L-1}$ for some $\lambda \in \mathbb{C} \setminus \{0\}$. By Lemma \ref{lemm:homogendiff}, $\dot{\mathbf{w}}' = (\lambda_{L-2}w_{L-2}, \ldots, \lambda_0w_0)$ for some $\lambda_i \in \mathbb{C} \setminus \{0\}$ such that $\lambda_{L-2} + r \lambda_{r-3} + \cdots + r^{L-2} \lambda_0 = \lambda$. But then
% \begin{align}
% \dot{\mathbf{w}} &= (-r \lambda w_{L-1}, \lambda_{L-2}w_{L-2}, \ldots, \lambda_0 w_0) = \\
% &= \lambda (-rw_{L-1}, w_{L-2}, 0, \ldots, 0) + (0, (\lambda_{L-2} - \lambda)w_{L-2}, \lambda_{L-3}w_{L-3}, \ldots, \lambda_{0} w_0).
% \end{align}
% Since $(\lambda_{L-2} - \lambda) + r \lambda_{L-3} + \cdots + r^{L-2}\lambda_0 = 0$, it follows from Lemma \ref{lemm:homogendiff} that the second summand of the above expression belongs to the kernel of $ \textnormal{J}_\mathbf{w'} \varphi$, and the desired result follows from the inductive hypothesis.  

%\textit{Case 2:} $k_{L-1} > 1$. 
  



% \section{THE DEGREE OF THE NEUROMANIFOLD}

% The degree of a projective variety quantifies the number of intersection points between the variety and a general linear space of complementary dimension. In other words, the degree provides a measure of how the variety is embedded in projective space, reflecting its geometric complexity. Given a projective variety \( \mathcal{V} \subset \mathbb{P}^n \), we denote its degree by \( \deg(\mathcal{V}) \). According to \cite{draisma2016euclidean}, it is known that \( \edd(\mathcal{V}) \geq \deg(\mathcal{V}) \).

% For the neuromanifold, as shown in \cite[Proposition 6.11]{kozhasov2023minimal}, the degree can be computed as follows:
% \begin{equation}
%     \label{eq:deg_neuro}
%     \deg(\mathbb{P}\mathcal{M}_{\mathbf{d}, \mathbf{k}, \mathbf{s}, r}) = (|\mathbf{k}|-L)!\prod_{0 \leq j < L} \frac{r^{(L - j - 1)(k_j - 1)}}{(k_j - 1)!}.
% \end{equation}
% From Equation \ref{eq:deg_neuro}, we observe that the degree of the neuromanifold exhibits exponential growth when the filter sizes \( k_j \) are close in value or when the parameter \( r \) increases -- see Figure \ref{fig:neur}.

% \begin{table}[!h]
%     \centering
%     \renewcommand{\arraystretch}{1.5}
%     \setlength{\tabcolsep}{12pt} % Adjust column width
%     \begin{tabular}{c|ccccc}
%         \diagbox[width=6em, height=2em]{\footnotesize $r$}{\footnotesize $t$} & \footnotesize $2$ & \footnotesize $3$ & \footnotesize $4$ & \footnotesize $5$ & \footnotesize $6$ \\
%         \hline
%         \footnotesize $1$ & 6 & 39 & 284 & 2205 & 17730 \\
%         \footnotesize $2$ & 14 & 219 & 3772 & 68405 & 1277898 \\
%         \footnotesize $3$ & 22 & 543 & 14684 & 417005 & 12186066 \\
%         \footnotesize $4$ & 30 & 1011 & 37244 & 1439205 & 57202074 \\
%         \footnotesize $5$ & 38 & 1623 & 75676 & 3699005 & 185917794 \\
%         \footnotesize $6$ & 46 & 2379 & 134204 & 7933205 & 482134890 \\
%     \end{tabular}
%     \caption{The generic Euclidean distance degree of neuromanifolds with $L=2$ and $t=k_1=k_2$.}
%     \label{tab:your_label}
% \end{table}



 
% Therefore, from Theorem \ref{th:segverdeg} it follows that:
% \begin{align}
%     &\edd(\overline{\mathcal{M}}_{\mathbf{d}, \mathbf{k}, \mathbf{s}}) = \edd(V_{\mathbf{m}, \mathbf{d}})= \\
%     \label{eq:ed_seg}
%     &= \sum_{0 \leq i \leq |\mathbf{d}|} (-1)^i \left( 2^{|\mathbf{d}| + 1 - i} - 1 \right)(|\mathbf{d}| - i)! \sum_{\substack{|\mathbf{\alpha}|=i \\ \forall j \ \alpha_j < k_j}} \prod_{0 \leq j < L} \binom{k_j}{\alpha_j} (k_j - \alpha_j - 1)! \ r^{(L - j -1)(k_j - \alpha_j - 1)},
% \end{align}
% where $|\mathbf{d}| = \sum_{0 \leq j < L}k_j -L$. The asymptotic behavior \vs{In what sense? when the degree of Veronese goes up or filter sizes?} of the above formula can be expressed as \vs{Is this straight forward?}: 
% \begin{align}
%     &\mathcal{O}\left(2^{|\mathbf{d}|} \ |\mathbf{d}|! \  r^{\sum_{j}(L-j-1)(k_j - 1)} \prod_{0 \leq j < L
%     }k_j ! \right)  = \\
%      = &\mathcal{O}\left( 2^{|\mathbf{d}|} \ |\mathbf{d}|! \  r^{(L-1)|\mathbf{d}|} \ \prod_{0 \leq j < L}k_j !  \right).
% \end{align}
% \vs{Degree and gEDdegree are two important invariants attached to an algebraic variety.
% They both can be written in terms of $k_j$ and $r$. The expectation is that these invariants increase as $k_j$ or $r$ increase. This is nice as the degree shows how twisted the neuromanifold is!
% I will compute the degree of Segre-Veronesne variety. Also, some figures might show that the neuromanifold acts like a space-filling manifold as $r$ or $k_j$ increases!   }

% \subsection{Complexity of the Neuromanifold}
% The computation of the $\edd$ of a projective variety can be done by summing its polar degrees. Specifically, the $\ell$-th polar degree of the Segre-Veronese variety $V_{\mathbf m,\mathbf d}$, for $\ell \in {0, \ldots, |\mathbf d|}$, is given by:
% \begin{equation}
%         \delta_\ell (V_{\mathbf m,\mathbf d}) = \sum_{i=0}^{|\mathbf d| - \ell} (-1)^i \binom{|\mathbf d| + 1 - i}{\ell + 1} (|\mathbf d| - i)!  
%       \sum_{\substack{|\mathbf{\alpha}|=i \\ \forall j \ \alpha_j < k_j}} \prod_{0 \leq j < L} \binom{k_j}{\alpha_j} (k_j - \alpha_j - 1)! \ r^{(L - j -1)(k_j - \alpha_j - 1)}.
% \end{equation}
% The degree of the variety $\deg(V_{\mathbf m,\mathbf d})$ is equal to $\delta_{|\mathbf d|}(V_{\mathbf m,\mathbf d})$ and thus can be computed as follows:
% \[
%  |\mathbf d|!  \prod_{0 \leq j < L}  (k_j  - 1)! \ r^{(L - j -1)(k_j - 1)}.
% \]
% By the above formula, we can see that an increase in $k_j$ or $r$ will result in a higher degree of the neuromanifold. In other words, the neuromanifold will have a larger volume with respect to the Fubini-Study metric.
% Moreover, by \parencite[Chapter 1, Corollary 5.11]{gelfand1994discriminants}, when $r > 1$, the dual variety $V^{\vee}_{\mathbf m,\mathbf d}$ is a hypersurface, and its degree is:
% \[
% \deg(V^{\vee}_{\mathbf m,\mathbf d}) = \delta_0(V_{\mathbf m,\mathbf d}) = \sum_{i=0}^{|\mathbf d|} (-1)^i (|\mathbf d| + 1 - i)! \sum_{\substack{|\mathbf{\alpha}|=i \\ \forall j \ \alpha_j < k_j}} \prod_{0 \leq j < L} \binom{k_j}{\alpha_j} (k_j - \alpha_j - 1)! \ r^{(L - j -1)(k_j - \alpha_j - 1)}.
% \]
%  \vs{I will provide further elaboration within this section. Additionally, a crucial point I will include is: An affine section of the neuromanifold exhibits a notably reduced algebraic complexity while maintaining its original degree! This is useful for overparamaterized convolutional network.}
% \subsection{Optimization}

% \section{Euclidean Distance Optimization}
% Degree and Euclidean distance degree are two fundamental invariants for a given variety \(\mathcal{V}\).

% The \emph{degree} of \(\mathcal{V}\), denoted by \(\deg(\mathcal{V})\), measures its intersection complexity—that is, the number of points at which \(\mathcal{V}\) intersects a generic linear space of complementary dimension. For instance, for a curve in the plane, it corresponds to the number of intersections with a generic line.

% On the other hand, \emph{Euclidean distance degree} of $\mathcal{V}$, denoted by $\operatorname{EDdeg}(\mathcal{V})$, aims to count the number of all critical points of the following quadratic distance function
% \begin{equation}
% \label{eq:dist}
%     \operatorname{dist_u: \mathcal{V}_{\operatorname{reg}} \to \mathbb{R}}; \quad v \mapsto \| v-u \|^2,
% \end{equation}
% where $\mathcal{V}_{\operatorname{reg}}$ is the smooth locus of $\mathcal{V}$ and $u\in \mathbb{R}^n$ is a generic data. Geometrically, this means
% \begin{equation}
% \label{eq:perp_tangent}
%     v\in \mathcal{V}_{\operatorname{reg}}; \quad v-u \perp \mathcal{T}_v \mathcal{V},
% \end{equation}
%  where $\mathcal{T}_v \mathcal{V}$ is the tangent space of $\mathcal{V}$ at the point $v$.  In algebraic terms, \eqref{eq:perp_tangent} can be expressed as a system of polynomial equations in the coordinates of \(v\) and \(u\). According to \parencite{draisma2016euclidean}, for almost all data points \(u \in \mathbb{R}^n\), this system of polynomials has a fixed number of solutions (including complex solutions), and this number is referred to as \(\operatorname{EDdeg}(\mathcal{V})\).
% \begin{exmp}
%     Consider a unit circle described by the equation \( x_1^2 + x_2^2 - 1 = 0 \). At a given point \( v = (x_1, x_2) \) on the circle, the tangent space is described by the span of the vector \( (x_2, -x_1) \). For a given data point \( u = (u_1, u_2) \in \mathbb{R}^2 \), the corresponding polynomial system for \eqref{eq:perp_tangent} is
%     \begin{equation}
%     \label{eq:circ_equ}
%     \begin{cases}
%         u_2 x_1 - u_1 x_2 &= 0 \\
%         x_1^2 + x_2^2 - 1 &= 0
%     \end{cases}
%     \end{equation}
% It is easy to verify that \eqref{eq:circ_equ} has only two solutions, both of which are real, whenever \( u \neq (0,0) \). Thus, the Euclidean distance degree of a circle is \( 2 \).
% \end{exmp}
% The isotropic quadric that defines the squared distance in \eqref{eq:dist} is very special. In our analysis, we are interested in all critical points encountered in
% \begin{equation}
% \label{eq:dist}
%     \operatorname{dist}_{A,u}: \mathcal{V}_{\operatorname{reg}} \to \mathbb{R}; \quad v \mapsto \| v-u \|_A^2=(v-u)^\top A (v-u),
% \end{equation}
% where \( A \) is a generic symmetric positive definite matrix. This number is known as the generic Euclidean distance degree of \( \mathcal{V} \), denoted by \( \operatorname{gEDdeg}(\mathcal{V}) \). It is known that for any variety \( \mathcal{V} \), \( \operatorname{gEDdeg}(\mathcal{V}) \ge \operatorname{EDdeg}(\mathcal{V}) \) (cf.~\parencite{maxim2020defect}).

% We now analyze the generic Euclidean distance degree of Segre--Veronese varieties. 
% \begin{thm}[\parencite{kozhasov2023minimal}]\label{th:segverdeg}
% The generic Euclidean Distance degree of the Segre-Veronese variety is
% \begin{equation}
% \label{eq:ged_segve_chern}
%     \edd(V_{\mathbf{m}, \mathbf{d}}) = \sum_{0 \leq i \leq |\mathbf{d}|} (-1)^i \left( 2^{|\mathbf{d}| + 1 - i} - 1 \right)(|\mathbf{d}| - i)! \sum_{\substack{|\mathbf{\alpha}|=i \\ \forall j \ \alpha_j \leq d_j}} \prod_{1\leq j \leq k} \binom{d_j + 1}{\alpha_j} (d_j - \alpha_j)! \ m_j^{d_j - \alpha_j},
%     \end{equation}

%     where $|\mathbf{d}| = d_1 + \cdots + d_k$. 
% \end{thm}
% The formula \eqref{eq:ged_segve_chern}  can be rewritten to the following way:
% \begin{equation}
%     \sum_{\ell=0}^{|\mathbf{d}|} \sum_{i=0}^{|\mathbf d| - \ell} (-1)^i \binom{|\mathbf d| + 1 - i}{\ell + 1} (|\mathbf d| - i)! \sum_{\substack{|\mathbf{\alpha}|=i \\ \forall j \ \alpha_j < d_j}} \prod_{0 \leq j < L} \binom{d_j+1}{\alpha_j} (d_j - \alpha_j)! \ r^{(L - j -1)(d_j - \alpha_j)}.
% \end{equation}
% \subsection{Rephrasing Learning as Distance Minimization}
% We now formulate the optimization problem in a convolutional polynomial network as a nearest point problem. Given input data \( X \subset \mathbb{R}^{d_0 \times N} \) and corresponding labels \( Y \subset \mathbb{R}^{d_L \times N} \), the task is to minimize the following quadratic loss:
% \[
% \mathcal{L}(w) = \| \varphi_w(X) - Y \|^2 = \sum_{(x, y) \in X \times Y} \| \varphi_w(x) - y \|^2.
% \]
% Over the neuromanifold and its Zariski closure, this optimization problem translates to:
% \begin{equation}
%     \min_{w \in \mathcal{M}_{\mathbf{d}, \mathbf{k}, \mathbf{s}}} \| \varphi_w(X) - Y \|^2 \quad \text{and} \quad \min_{w \in \overline{\mathcal{M}}_{\mathbf{d}, \mathbf{k}, \mathbf{s}}} \| \varphi_w(X) - Y \|^2,
%     \label{eq:opt_zariski_closure}
% \end{equation}
% respectively. A matrix representation \( \widetilde{W} \) of \( \varphi_w \) has a convolutional format with filter \( \widetilde{w} \) of size \( \widetilde{k} \) and stride \( \widetilde{s} \). By \parencite{shahverdi2024algebraic}[Equation 4.4], for sufficiently large \( N \), \eqref{eq:opt_zariski_closure} translates to:
% \begin{equation}
%     \label{eq:opt_trace}
%     \min_{\widetilde{W} \in \overline{\mathcal{M}}_{\mathbf{d}, \mathbf{k}, \mathbf{s}}} \operatorname{tr}[(\widetilde{W} - U)\widetilde{X}\widetilde{X}^\top (\widetilde{W} - U)^\top],
% \end{equation}
% where \( U \) is a convolutional matrix with filter \( u \) of size \( \widetilde{k} \) and stride \( \widetilde{s} \), obtained by projecting an arbitrary matrix of the same size as \( \widetilde{W} \) onto the linear space of convolutional matrices of that size. \( \widetilde{X} \) is a matrix of size \( \widetilde{d}_0 \times \widetilde{d}_0 \), corresponding to the input matrix \( X \) after passing through the Segre--Veronese map.

% Furthermore, by \parencite{shahverdi2024algebraic}[Equations 4.5 and 4.8], there exists a linear map \( \psi \) such that:
% \begin{align}
%     \label{eq:psi}
%     \psi: \mathbb{R}^{\widetilde{d}_0 \times \widetilde{d}_0} \to \mathbb{R}^{\widetilde{k} \times \widetilde{k}}, \quad M \mapsto \left[ \sum_{m=0}^{\widetilde{d}_L-1} M_{i+\widetilde{s}m, j+\widetilde{s}m} \right]_{ij},
% \end{align}
% which converts \eqref{eq:opt_trace} into an optimization over the filters as follows:
% \begin{equation}
%     \label{eq:opt_trace_filter}
%     \min_{\widetilde{w} \in \overline{\mathcal{M}}_{\mathbf{d}, \mathbf{k}, \mathbf{s}}} \operatorname{tr}[(\widetilde{w} - u)\psi(\widetilde{X}\widetilde{X}^\top) (\widetilde{w} - u)^\top].
% \end{equation}
% Note that the Segre--Veronese variety does not lie in any vector space of positive codimension; hence \( \psi(\widetilde{X}\widetilde{X}^\top) \) is a generic matrix of size \( \widetilde{k} \times \widetilde{k} \). Therefore, the number of critical points in \eqref{eq:opt_trace_filter} is the generic Euclidean distance degree of \( \overline{\mathcal{M}}_{\mathbf{d}, \mathbf{k}, \mathbf{s}} \), or equivalently, the corresponding Segre--Veronese variety.
% \iffalse
% Let us now discuss the optimization scenario arising when training a polynomial neural network on a regression problem. Specifically, following \parencite{kubjas2024geometry}, we will reformulate the classical squared error loss w.r.t. a dataset as a distance minimization objective over the neuromanifold. To this end, consider a finite dataset $\mathcal{D} \subseteq \mathbb{R}^n \times \mathbb{R}$ sampled from a homogeneous polynomial $p \in \symm{r^{L-1}}{n}$, i.e., $p(x) = y$ for all $(x,y) \in \mathcal{D}$. The objective that the network minimizes is defined for $w \in W $ as: 
% \begin{equation}\label{eq:loss}
% \mathcal{L}(w) = \sum_{(x,y) \in \mathcal{D}} (\varphi_w(x) - y )^2.
% \end{equation}
% % Suppose that $|\mathcal{D}|$ is large enough such that, for generic $\mathcal{D}$, there is a unique polynomial function $p_\mathcal{D}$ interpolating the dataset, i.e., $p_{\mathcal{D}}(x) = y$ for all $(x,y) \in \mathcal{D}$. 

% Then Equation \ref{eq:loss} can be rephrased quadratically w.r.t. $\varphi_w$:
% \begin{equation}
% \mathcal{L}(w) = (\varphi_w - p)^T E_{\mathcal{D}}(\varphi_w - p), \hspace{2em}  E_\mathcal{D} = \left( \sum_{(x,y) \in \mathcal{D}} x^{\alpha + \beta} \right)_{\alpha, \beta \in \binom{n+d-1}{d}}
% \end{equation}
% Note that for generic $\mathcal{D}$ with fixed large cardinality, $E_\mathcal{D}$ is a generic symmetric matrix. This reduces the problem of training a polynomial neural network with to the problem of finding the element of the neuromanifold which is closest to a generic element of the ambient space w.r.t. a generic quadratic form. The number of critical points in $\symm{d}{n}$ is therefore controlled by the (generic) Euclidean distance degree of $\mathcal{M}$. 
% \fi

% \newpage
% \bibliography{biblio}
% \bibliographystyle{alpha}

\end{document}
